\documentclass[11pt]{article}
\usepackage[margin=1.1in]{geometry}
\usepackage{amsmath,amssymb,amsthm,mathtools,booktabs,array,tabularx,longtable}
\usepackage[dvipsnames]{xcolor}
\usepackage[colorlinks=true,linkcolor=MidnightBlue,citecolor=MidnightBlue,urlcolor=MidnightBlue]{hyperref}
\usepackage{fancyhdr}
\pagestyle{fancy}\fancyhf{}\renewcommand{\headrulewidth}{0pt}
\fancyfoot[R]{\thepage}
\fancyfoot[L]{\scriptsize\color{gray}Generated by the Physics Graph, not the version of record. Official version: \href{https://doi.org/10.5281/zenodo.24402873}{doi:10.5281/zenodo.24402873}}
\fancypagestyle{plain}{\fancyhf{}\fancyfoot[R]{\thepage}\fancyfoot[L]{\scriptsize\color{gray}Generated by the Physics Graph, not the version of record. Official version: \href{https://doi.org/10.5281/zenodo.24402873}{doi:10.5281/zenodo.24402873}}\renewcommand{\headrulewidth}{0pt}}
\newtheorem{theorem}{Theorem}[section]
\newtheorem{proposition}[theorem]{Proposition}
\newtheorem{lemma}[theorem]{Lemma}
\newtheorem{corollary}[theorem]{Corollary}
\newtheorem{conjecture}[theorem]{Conjecture}
\newtheorem{refutation}[theorem]{Refutation}
\newtheorem{prediction}[theorem]{Prediction}
\newtheorem{principle}[theorem]{Principle}
\newtheorem{observation}[theorem]{Observation}
\newtheorem{axiom}[theorem]{Axiom}
\theoremstyle{definition}
\newtheorem{definition}[theorem]{Definition}
\newtheorem{fact}[theorem]{Fact}
\newtheorem{assembled}[theorem]{Assembly}
\newtheorem{postulate}[theorem]{Postulate}
\newtheorem{premise}[theorem]{Premise}
\newtheorem{hypothesis}[theorem]{Hypothesis}
\newtheorem{remark}[theorem]{Remark}
\newtheorem{example}[theorem]{Example}
\theoremstyle{plain}
\setlength{\marginparwidth}{0.8in}
\newcommand{\changedmark}[2]{\marginpar{\raggedright\scriptsize\color{#1}#2}}
\providecommand{\Aa}{}\renewcommand{\Aa}{\mathcal{A}}
\providecommand{\Ad}{}\renewcommand{\Ad}{\widetilde{\mathrm{Ad}}}
\providecommand{\Aev}{}\renewcommand{\Aev}{\mathcal{A}^{+}}
\providecommand{\Aod}{}\renewcommand{\Aod}{\mathcal{A}^{-}}
\providecommand{\B}{}\renewcommand{\B}{\mathrm{B}}
\providecommand{\C}{}\renewcommand{\C}{\mathbb{C}}
\providecommand{\CAR}{}\renewcommand{\CAR}{\mathrm{CAR}}
\providecommand{\CS}{}\renewcommand{\CS}{{\mathbb{S}^{+}}}
\providecommand{\CSm}{}\renewcommand{\CSm}{{\mathbb{S}^{-}}}
\providecommand{\Cl}{}\renewcommand{\Cl}{\mathrm{Cl}}
\providecommand{\Con}{}\renewcommand{\Con}{\mathcal{C}}
\providecommand{\Dl}{}\renewcommand{\Dl}{\Delta}
\providecommand{\EN}{}\renewcommand{\EN}{E_{N}}
\providecommand{\End}{}\renewcommand{\End}{\mathrm{End}}
\providecommand{\F}{}\renewcommand{\F}{\mathbb{F}}
\providecommand{\G}{}\renewcommand{\G}{\Gamma}
\providecommand{\HH}{}\renewcommand{\HH}{\mathbb{H}}
\providecommand{\IS}{}\renewcommand{\IS}{I\sigma_3}
\providecommand{\J}{}\renewcommand{\J}{\mathbf{J}}
\providecommand{\LH}{}\renewcommand{\LH}{L_{H}}
\providecommand{\Mat}{}\renewcommand{\Mat}{\mathrm{M}}
\providecommand{\Mo}{}\renewcommand{\Mo}{M^{\circ}}
\providecommand{\Mp}{}\renewcommand{\Mp}{M^{+}}
\providecommand{\Msun}{}\renewcommand{\Msun}{M_{\odot}}
\providecommand{\Mt}{}\renewcommand{\Mt}{\widetilde{M}}
\providecommand{\Ncone}{}\renewcommand{\Ncone}{\mathcal{N}}
\providecommand{\Neff}{}\renewcommand{\Neff}{N_{\mathrm{eff}}}
\providecommand{\PCS}{}\renewcommand{\PCS}{\mathbb{P}{\mathbb{S}^{+}}}
\providecommand{\PT}{}\renewcommand{\PT}{\mathbb{PT}}
\providecommand{\Pin}{}\renewcommand{\Pin}{\mathrm{Pin}}
\providecommand{\Q}{}\renewcommand{\Q}{\mathbb{Q}}
\providecommand{\R}{}\renewcommand{\R}{\mathbb{R}}
\providecommand{\Ric}{}\renewcommand{\Ric}{\mathrm{Ric}}
\providecommand{\SBH}{}\renewcommand{\SBH}{S_{\mathrm{BH}}}
\providecommand{\Sc}{}\renewcommand{\Sc}{\mathbb{S}}
\providecommand{\Sch}{}\renewcommand{\Sch}{\mathrm{P}}
\providecommand{\Sig}{}\renewcommand{\Sig}{\Sigma}
\providecommand{\Som}{}\renewcommand{\Som}{S_\omega}
\providecommand{\Sp}{}\renewcommand{\Sp}{\mathbb{S}}
\providecommand{\Sph}{}\renewcommand{\Sph}{\mathrm{S}}
\providecommand{\Spi}{}\renewcommand{\Spi}{S_\pi}
\providecommand{\Spin}{}\renewcommand{\Spin}{\mathrm{Spin}}
\providecommand{\Sth}{}\renewcommand{\Sth}{S_{\mathrm{th}}}
\providecommand{\T}{}\renewcommand{\T}{\mathbb{T}}
\providecommand{\TH}{}\renewcommand{\TH}{T_{\mathrm{H}}}
\providecommand{\Tr}{}\renewcommand{\Tr}{\operatorname{Tr}}
\providecommand{\UHF}{}\renewcommand{\UHF}{\mathrm{UHF}}
\providecommand{\Z}{}\renewcommand{\Z}{\mathbb{Z}}
\providecommand{\Zc}{}\renewcommand{\Zc}{\mathrm{Z}}
\providecommand{\cmg}{}\renewcommand{\cmg}{\mathrm{cm^{2}\,g^{-1}}}
\providecommand{\dS}{}\renewcommand{\dS}{\mathrm{dS}}
\providecommand{\dsc}{}\renewcommand{\dsc}{\mathrm{dsc}}
\providecommand{\dtPG}{}\renewcommand{\dtPG}{dt}
\providecommand{\dtS}{}\renewcommand{\dtS}{dt_{\mathrm{S}}}
\providecommand{\ellv}{}\renewcommand{\ellv}{\ell}
\providecommand{\er}{}\renewcommand{\er}{e_r}
\providecommand{\et}{}\renewcommand{\et}{e_t}
\providecommand{\ev}{}\renewcommand{\ev}{\mathrm{ev}}
\providecommand{\gl}{}\renewcommand{\gl}{\mathfrak{gl}}
\providecommand{\hatV}{}\renewcommand{\hatV}{\widehat{V}}
\providecommand{\hb}{}\renewcommand{\hb}{\bar{h}}
\providecommand{\hf}{}\renewcommand{\hf}{\hat{f}}
\providecommand{\id}{}\renewcommand{\id}{\mathrm{id}}
\providecommand{\kpc}{}\renewcommand{\kpc}{\mathrm{kpc}}
\providecommand{\lP}{}\renewcommand{\lP}{\ell_{P}}
\providecommand{\mD}{}\renewcommand{\mD}{\mathbf{m}_{D}}
\providecommand{\mb}{}\renewcommand{\mb}{\bar m}
\providecommand{\mnu}{}\renewcommand{\mnu}{\Sigma m_{\nu}}
\providecommand{\nInf}{}\renewcommand{\nInf}{n_\infty}
\providecommand{\nO}{}\renewcommand{\nO}{n_0}
\providecommand{\ninf}{}\renewcommand{\ninf}{n_{\infty}}
\providecommand{\nz}{}\renewcommand{\nz}{n_{0}}
\providecommand{\proj}{}\renewcommand{\proj}{\operatorname{proj}}
\providecommand{\pv}{}\renewcommand{\pv}{\pi}
\providecommand{\rad}{}\renewcommand{\rad}{\mathrm{rad}}
\providecommand{\sgn}{}\renewcommand{\sgn}{\operatorname{sgn}}
\providecommand{\shA}{}\renewcommand{\shA}{\mathcal{A}}
\providecommand{\shC}{}\renewcommand{\shC}{\mathcal{C}}
\providecommand{\shE}{}\renewcommand{\shE}{\mathcal{E}}
\providecommand{\shJ}{}\renewcommand{\shJ}{\mathcal{J}}
\providecommand{\shK}{}\renewcommand{\shK}{\mathcal{K}}
\providecommand{\shL}{}\renewcommand{\shL}{L}
\providecommand{\shW}{}\renewcommand{\shW}{\mathcal{W}}
\providecommand{\sll}{}\renewcommand{\sll}{\mathfrak{sl}}
\providecommand{\so}{}\renewcommand{\so}{\mathfrak{so}}
\providecommand{\su}{}\renewcommand{\su}{\mathfrak{su}}
\providecommand{\supp}{}\renewcommand{\supp}{\operatorname{supp}}
\providecommand{\tr}{}\renewcommand{\tr}{\operatorname{tr}}
\providecommand{\uu}{}\renewcommand{\uu}{\mathfrak{u}}
\providecommand{\br}{}\renewcommand{\br}[1]{\langle #1\rangle}
\providecommand{\bra}{}\renewcommand{\bra}[1]{\langle #1|}
\providecommand{\file}{}\renewcommand{\file}[1]{\texttt{#1}}
\providecommand{\hp}{}\renewcommand{\hp}[2]{( #1 \,|\, #2 )}
\providecommand{\ip}{}\renewcommand{\ip}[2]{\langle #1, #2\rangle}
\providecommand{\ket}{}\renewcommand{\ket}[1]{|#1\rangle}
\providecommand{\rel}{}\renewcommand{\rel}[2]{D\bigl(#1 \,\|\, #2\bigr)}
\providecommand{\rev}{}\renewcommand{\rev}[1]{\widetilde{#1}}
\providecommand{\slashed}{}\renewcommand{\slashed}[1]{{#1}\mkern-10mu/}
\title{\bfseries Geometry First\\[4pt]\large\itshape From the Bit to the Numbers}
\author{Leon Butler\\ \normalsize\href{mailto:lbutler.math@gmail.com}{\texttt{lbutler.math@gmail.com}}\\ \normalsize ORCID: \href{https://orcid.org/0009-0006-4449-5617}{0009-0006-4449-5617}}
\date{Generated by the Physics Graph, 20 November 2026 \textperiodcentered{} Transport Papers, T0a\\[3pt]\small Official version: \href{https://doi.org/10.5281/zenodo.24402873}{doi:10.5281/zenodo.24402873}}
\begin{document}
\maketitle
\begin{abstract}
\normalsize
The primer of the Transport Papers, like the textbooks, builds geometric algebra on the real numbers: numbers first, geometry on top. This companion note runs the construction in the opposite order. It starts from a single distinction with two states and asks how much of the algebra, and how many of the numbers, follow from it. From one bit of non-commutativity and integer superposition it proves the one rule $v^{2} = Q(v)$ as a theorem and obtains the integer Clifford algebras, the signature appearing as a quadratic form over the bit. Demanding that products be undone forces the rationals. After completion, which supplies the real numbers as scalars, it proves Sylvester's law of inertia, that every nonconstant real polynomial has its roots in $\Cl(0,1)$, where a direction of square $-1$ stands in for the imaginary unit, and Frobenius' theorem, so that the complex numbers and the quaternions appear as members of the family the bit produces. Setting and clearing a bit are shown to be exactly a hyperbolic pair over every commutative ring, the relations of one fermionic mode, while the integer algebra of the ladder contains no such switch. A closing section shows what tests leave of a register of switches: the pinching onto the functions of the configuration, which is the descent of T5. Four premises and one demand are chosen, and every other step is proved. The ingredients are classical and attributed, and none of it is evidence for the transport principle.
\end{abstract}
\medskip
\noindent\textbf{MSC 2020.} Primary 15A66; Secondary 11E04, 16K20.\\
\textbf{Keywords.} geometric algebra; Clifford algebra; the bit; integer Clifford algebras; Sylvester's law of inertia; Frobenius theorem; fermionic mode.
\newpage
\noindent \textbf{Authorship and disclosure. }The founding idea of the series is the author's: that division by zero in a Clifford algebra should be read not as a prohibition but as a transport into a higher-dimensional algebra, and that the transported division need not be accompanied by an inverse operation. The direction of investigation and the critical review of the results are likewise the author's. The mathematical development and the drafting of the text were carried out in extended collaboration with Claude Fable 5 (Anthropic), an AI system.

\noindent\fbox{\parbox{\dimexpr\linewidth-2\fboxsep-2\fboxrule}{\textbf{This is not the version of record.} The official version of this paper is the one deposited on Zenodo, \href{https://doi.org/10.5281/zenodo.24402873}{doi:10.5281/zenodo.24402873}. This PDF was generated by the Physics Graph on 20 November 2026 and assembled from approved parts of its store: the results, proofs and text of the paper as the store holds them, each result with its premises and machine checks. It is regenerated when a result it contains is corrected, so it may differ from the deposited paper. Since the deposited version 1.0: 13 results changed, 12 cited results added and 1 deposited result withdrawn. Each revised or added result is marked in the margin (\textcolor{BrickRed}{corrected}, \textcolor{BurntOrange}{changed}, \textcolor{gray}{reworded}), and the last section, \emph{Changes since the deposited version}, lists them with the reason for each.}}
\newpage

\section{Geometry before number}

\phantomsection\label{c:note.geometry-first.setting-geometry-before-number}
The author's view (and likely many people's view) is that the universe is not calculating with numbers. Nothing in nature looks up a real number to decide what happens next. Whatever is happening is more basic than arithmetic, and the numbers we use to describe it are a compact summary of it: extremely good for calculation, but not the process itself. Two older voices point the same way. Galileo wrote that the book of nature is written in the language of mathematics, its characters being triangles, circles and other geometric figures \cite{Galileo1623}. Wheeler proposed that every physical quantity draws its meaning from yes-or-no questions: ``it from bit'' \cite{Wheeler1990}. Between them they name the two things this note puts first, geometry and the bit, although Wheeler went further than this note does and rejected the continuum that its last premise posits. Built out from the bit, the numbers and the algebra lead to geometric algebra. The premises were chosen so that this is where they lead, which limits what the derivation can show (Remark~\emph{Argued: geometric algebra as the algebra of physics} [T0a, Rem 7.6]); the author reads it as further support for the view, argued most fully by Hestenes \cite{Hestenes2003Oersted}, that geometric algebra is \emph{the} mathematical framework for the natural sciences.

Mathematics is usually built the other way round. The primer T0 \cite{ButlerT0} begins where the textbooks begin, with the real numbers, a vector space over them and a quadratic form: numbers come first and geometry is built on top. This note reverses the order. It starts from a single distinction, something or nothing, and asks how much of the algebra, and how many of the numbers, can be made to follow from it. Geometry here means the sign and the signature, a quadratic form over the bit, which appear before any coefficient ring does.

The answer is: almost all of it. Four premises and one demand are chosen, and everything else is proved. The sign and the signature reduce to one posited element and one free choice. The one rule $v^{2} = Q(v)$, which the primer takes as its starting point \cite[Section 2]{ButlerT0}, becomes a theorem. The integers enter as counts, and the demand that products can be undone forces the rationals. One further premise, completion, supplies the real numbers as the scalars; the complex numbers and the quaternions are then found as members of the family the bit produces, $\Cl(0,1)$ and $\Cl(0,2)$, rather than as its foundation. Alongside this ladder runs a second thread, what can be done to a single bit, which leads to the switch of the section ``The switch: what a bit does'' and to what its tests leave in the section ``What the tests leave of a switch''.

The starting point has a history. Spencer-Brown's calculus of indications begins from a distinction, and his law of crossing is the rule of (F1) that a distinction made twice returns to nothing \cite{SpencerBrown1969}. Kauffman has connected such distinctions, through his iterants, to Clifford algebras and to fermions \cite{Kauffman2021}. Von Weizs\"acker took binary alternatives, his ur-alternatives, as the foundation of physics \cite{Weizsacker1985}. And a Clifford algebra is known to be a group algebra twisted by a sign \cite{AlbuquerqueMajid2002}. What this note adds is the ladder itself, in exact integer arithmetic from each premise to the next, the switch theorem with its failure over the integers, and the pinching that joins the two.

\phantomsection\label{c:note.geometry-first.setting-geometry-before-number-2}
\textbf{Status.} None of this proves the view the section opened with. The mathematics shows how much follows once the premises are granted, and no more; whether that says anything about the world is a reading, set out with its objections in Remark~\emph{Argued: geometric algebra as the algebra of physics} [T0a, Rem 7.6]. To keep the two apart, every result of the note carries one of three grades, and a sentence that only interprets a result says that it is a reading.
\begin{center}
\begin{tabular}{@{}lp{0.72\textwidth}@{}}
\toprule
\emph{chosen} & a premise of the foundation, (F1) to (F4), or the demand (D): assumed, not derived \\
\emph{proved} & a theorem, proposition or lemma, with its proof given here \\
\emph{argued} & an inference from evidence, stated as a remark or marked as a reading, and never used in a proof \\
\bottomrule
\end{tabular}
\end{center}
The premises are mathematical: they posit structures and assert nothing about the world. Unlike the primer, which is expository and claims nothing new, the note's arrangement of the steps is the series' own, although each ingredient is classical and is attributed where it is used. Nothing in the primer, and nothing in the toolkit T1 imports from it, depends on this note.

\phantomsection\label{c:note.geometry-first.setting-geometry-before-number-3}
\textbf{How to read this note.} a figure of the original is the whole route. the section ``The bit and its repetition'' to~the section ``What the tests leave of a switch'' take it one step at a time: each states the step, says in a sentence how the proof goes, and works small cases by hand for most of its results. Every proof is in the section ``Proofs'', so a first reading can skip them all; the main text does not depend on them. Every worked example is machine-checked (the section ``Machine verification'').

\section*{What we assume}
The results below are taken as given in this paper; each result that rests on one says so.
\begin{postulate}[Premise (F4): the coefficients are completed to a complete ordered field]\label{c:synthesis.continuum-completion-premise}
\changedmark{BurntOrange}{changed since the deposit}
\emph{Premise (F4), completion.} The coefficients of the Clifford algebras, the field $\Q$ that division reaches from the integers, are extended to a \emph{complete ordered field}: a field $\R$ with a total order $<$ such that $a < b$ implies $a + c < b + c$, $a > 0$ and $b > 0$ imply $ab > 0$, and every nonempty subset of $\R$ that is bounded above has a least upper bound in $\R$. The rational Clifford algebras are extended to $\R$ by scalars, and $\R$ is the field over which the store's quadratic spaces (\emph{Real quadratic space} [algebra core]) and Clifford algebras (\emph{The Clifford algebra of a quadratic space} [algebra core]) are defined.

\emph{What is assumed.} Only that such a field exists and is taken as the coefficients. The inclusion $\Q \subseteq \R$ is not a further assumption: in an ordered field $1 = 1^{2} > 0$, so $m \cdot 1 > 0$ for every integer $m \ge 1$, the field has characteristic $0$, and $a/m \mapsto (a \cdot 1)(m \cdot 1)^{-1}$ is the unique embedding of $\Q$, order-preserving. The premise is a mathematical primitive of the series' foundation, not a statement about the physical world.

\emph{Status and cost.} Premise. It is the one step of the foundation that is not algebraic: the integer and rational algebras are countable, and the premise adds the limits. What it is needed for, and what it is not: square roots of positive elements, which diagonalise every quadratic form to squares $1$, $-1$, $0$, follow from it, but need only a field in which every positive element is a square (a Euclidean field), not completeness; roots of real polynomials of odd degree, and with them Frobenius' classification of the real division algebras, are obtained here from completeness, but hold in any real closed field (for example the real algebraic numbers); completeness proper is needed for limits such as the exponential of a bivector and the angle $\pi$ it defines.

\emph{What would retire it.} Since the premise is mathematical, no experiment refutes it; it would become unnecessary for a result that is rederived over a countable ordered field (for example a Euclidean or a real closed field of algebraic numbers), and it is unnecessary for every result whose statement and proof stay over $\Q$.
\end{postulate}

\begin{postulate}[Premise (F2): minimal non-commutativity]\label{c:synthesis.premise-minimal-non-commutativity}
\changedmark{BurntOrange}{changed since the deposit}
\emph{Premise (F2), minimal non-commutativity.} Grant Premise (F1), the bit, with $n \ge 1$ distinctions. A group $G$, generated by the directions $e_{1}, \dots, e_{n}$ and a bit $c$, extends the configurations $\mathbb{F}_{2}^{n}$ by that one central bit, and distinct directions do not commute:
\begin{enumerate}
\item[(F2a)] $c$ is central and $c^{2} = 1$;
\item[(F2b)] $G/\langle c\rangle$ is elementary abelian of order $2^{n}$ with basis the images of the $e_{i}$, so that $S \mapsto e_{S}\langle c\rangle$, $e_{S}$ the increasing product of the $e_{i}$ with $i \in S$, identifies $\mathbb{F}_{2}^{n}$ under symmetric difference with $G/\langle c\rangle$;
\item[(F2c)] $c \ne 1$ and $e_{i}e_{j} = c\,e_{j}e_{i}$ for all $i \ne j$.
\end{enumerate}
Which directions square to $c$ rather than to $1$ (each $e_{i}^{2}$ lies in $\langle c\rangle$ by (F2b)) is not fixed by the premise: every choice, the signature, is allowed and none is singled out.

\emph{Status and cost.} A chosen primitive of the foundation, mathematical and not physical. It is minimal in two senses: the extension is by the smallest nontrivial group, one bit, and the commutator of any two distinct directions is that bit. The clause $c^{2} = 1$ is forced by (F2b) and (F2c) when $n \ge 2$, and is independent of them when $n = 1$. Once $c^{2} = 1$ holds, centrality is implied as well, since $\langle c\rangle$ is then a normal subgroup of order $2$; so for $n \ge 2$ all of (F2a) follows from (F2b) and (F2c), and (F2a) has content only at $n = 1$, through its order clause (without it, a dihedral group with $c$ a rotation satisfies (F2b) and (F2c) with $c$ not central). What it costs: order now matters, by exactly one bit. What it does not say: anything about addition or coefficients (that is the premise of superposition), or which signature nature uses. No observation bears on it directly; the reading that the sign of a spinor under a full turn is this bit is argued, not proved.
\end{postulate}

\begin{postulate}[Premise (F3): superposition over the integers]\label{c:synthesis.premise-superposition-over-the-integers}
\changedmark{BurntOrange}{changed since the deposit}
\emph{Premise (F3), superposition over the integers.} Grant Premise (F2), with its group $G$ and central bit $c$. Signed configurations (the elements of $G$) can be added, with integer coefficients, and the sign bit acts as $-1$: the algebra of the foundation is
\[
A = \Z[G]/(c + 1),
\]
the integral group ring of $G$ modulo the two-sided ideal generated by $c + 1$.

\emph{Status and cost.} A chosen primitive of the foundation, mathematical and not physical. The coefficient ring is the integers, the ring the unit generates under addition; the rationals and the reals are not assumed here (the first is reached by dividing by the squares of directions, the second needs the premise of completion). Identifying $c$ with $-1$ is part of the choice: in the group ring $c$ is a central element with $c^{2} = 1$, and the other possibility, $c$ acting as $+1$ (the only other integer of square $1$), gives the commutative group ring of the configurations, in which order is forgotten. What it does not say: that the coefficients are a field, or anything about limits. No observation bears on it directly; that quantum superposition is this addition is a reading, argued and not proved.
\end{postulate}

\begin{postulate}[Premise (F1): the bit]\label{c:synthesis.premise-the-bit}
\changedmark{BurntOrange}{changed since the deposit}
\emph{Premise (F1), the bit.} The foundation takes as primitive a single distinction with two states, written $0$ (nothing) and $1$ (something), combined by comparison: making a distinction and making it again returns to nothing, $1 + 1 = 0$. With the product $1 \cdot 1 = 1$ (both made) and $0$ otherwise, the states form the two-element field $\mathbb{F}_{2}$. Distinctions are independent: for $n \ge 0$ distinctions a configuration is the set $S \subseteq \{1, \dots, n\}$ of those made, and configurations combine by comparison, the symmetric difference $S \triangle T$. So the configurations form the vector space $\mathbb{F}_{2}^{n}$, with $2^{n}$ elements, $\binom{n}{k}$ of them containing $k$ distinctions.

\emph{Status and cost.} A chosen primitive of the foundation, mathematical and not physical: it posits a structure, the two-element field and its finite powers, and asserts nothing about the world. Nothing beneath it is available to derive it from. What it does not contain: a sign ($1 = -1$ in $\mathbb{F}_{2}$, so a distinction cannot be told from its reversal), an order between distinctions (comparison is commutative), and any magnitude or geometry; the sign and the order enter with the premise of minimal non-commutativity, superposition with the premise of superposition over the integers. No observation bears on it; what can fail is the reading that the algebra of physics is built on it, which is argued and not proved.
\end{postulate}

\section{The bit and its repetition}

\phantomsection\label{c:note.geometry-first.setting-bit-repetition}
No coefficient ring has been used: counting serves only to name the configurations. Counting the configurations, when numbers are later available, gives $2^{n}$ of them, of which $\binom{n}{k}$ have exactly $k$ distinctions on: the dimension count and the grade count of the primer \cite[Section 3]{ButlerT0}. In later language the configurations form the vector space $\F_{2}^{n}$ over the field with two elements, and that name is used below as a label only.

What the bit does not contain is as important. It has no sign, since a distinction cannot be told from its reversal. It has no order, since comparison is symmetric. It has no magnitude and no geometry.

\section{Order, and the arrival of the sign}

\phantomsection\label{c:note.geometry-first.setting-order-arrival-sign}
The second premise says only that the order in which two different distinctions are made can matter, and that when it does it matters by a single sign. Formally:

\phantomsection\label{c:note.geometry-first.setting-order-arrival-sign-2}
In words: the order in which two different distinctions are made now matters, and it matters by exactly one bit, the element $c$. This is the sense in which the non-commutativity is called minimal: every pair of different directions fails to commute by the same single element. The premise does not say which directions square to $c$ and which to $1$; every choice is allowed, and Theorem \ref{thm:bitsign}(v) shows that every choice occurs.

\phantomsection\label{c:note.geometry-first.setting-order-arrival-sign-3}
\noindent\emph{How the proof goes.} Any product of directions can be put in increasing order, each exchange of two different directions costing one $c$ and each repeated direction being replaced by its square; counting the costs gives $f$, and part (v) builds the group from that same rule. Example~\emph{Two bits} [T0a, Ex A] is the case $n = 2$. The proof is in Appendix~\ref{prf:thm:bitsign}.

So three things have appeared before any addition and before any coefficient ring: a sign (the element $c$, which will act as $-1$), a signature (the set $N$, which is a choice), and, as soon as $n \geq 2$, an element squaring to the sign whatever the signature: either some $e_{i}$ does, or two directions outside $N$ have $q(\{i,j\}) = 1$, so that $e_{i}e_{j}$ squares to $c$ (the unit bivector of \cite[Example 2.1]{ButlerT0}). The functions $q$ and $b$ are a quadratic form on the configurations and its polar form: the signature is a quadratic form over the bit.

\section{Superposition, and the one rule as a theorem}

\phantomsection\label{c:note.geometry-first.setting-superposition-one-rule-theorem}
The third premise lets elements be added, which brings in the integers as counts, and reads the sign as $-1$. Formally:

\phantomsection\label{c:note.geometry-first.setting-superposition-one-rule-theorem-2}
The integers enter here, as counts of how many times a configuration is added, and $-1$ enters as the image of $c$. For a commutative ring $R$ write $A_{R} = R[G]/(c+1)$ for the same construction with coefficients in $R$, and for $x \in R^{n}$ put
\[
Q(x) \;=\; \sum_{i \notin N} x_{i}^{2} \;-\; \sum_{i \in N} x_{i}^{2} .
\]

\begin{theorem}[The group of signed directions and its integer Clifford algebra]\label{c:synthesis.signed-directions-and-the-integer-clifford-algebra}
\changedmark{BurntOrange}{changed since the deposit}
Let $n \ge 1$. A \emph{group of signed directions} is a group $G$ with elements $e_{1}, \dots, e_{n}$ and $c$ such that
\begin{enumerate}
\item[(F2)] $c$ is central; $G/\langle c\rangle$ is elementary abelian of order $2^{n}$ with basis the images of $e_{1}, \dots, e_{n}$, so that $S \mapsto e_{S}\langle c\rangle$ identifies $\mathbb{F}_{2}^{n}$ (subsets of $\{1, \dots, n\}$ under symmetric difference $S \triangle T$) with $G/\langle c\rangle$, where $e_{S}$ is the increasing product of the $e_{i}$, $i \in S$; distinct directions do not commute, $c \ne 1$ and $e_{i}e_{j} = c\,e_{j}e_{i}$ for $i \ne j$; and $c^{2} = 1$.
\end{enumerate}
Then $e_{i}^{2} \in \langle c\rangle$; write $e_{i}^{2} = c^{s_{i}}$ with $s_{i} \in \{0, 1\}$, $N = \{i : s_{i} = 1\}$, $q = |N|$, $p = n - q$, and $Q(x) = \sum_{i \notin N} x_{i}^{2} - \sum_{i \in N} x_{i}^{2}$. Let
\begin{enumerate}
\item[(F3)] $A = \Z[G]/(c + 1)$, the integral group ring of $G$ modulo the two-sided ideal generated by $c + 1$.
\end{enumerate}
\begin{enumerate}
\item[(i)] \emph{(The sign.)} For $n \ge 2$ the clause $c^{2} = 1$ follows from the others; for $n = 1$ it is independent of them: the cyclic group of order $2m$ generated by $e_{1}$, with $c = e_{1}^{2}$, satisfies the other clauses for every $m \ge 2$ and has $c^{2} = 1$ only for $m = 2$. For every sign vector $(s_{1}, \dots, s_{n})$ a group of signed directions exists: $\{0, 1\} \times \mathbb{F}_{2}^{n}$ with $(a, S)(b, T) = (a + b + f(S,T) \bmod 2,\ S \triangle T)$, $f$ as in (ii), $c = (1, \emptyset)$, $e_{i} = (0, \{i\})$; it is unique up to an isomorphism fixing $c$ and the $e_{i}$. Every element of $G$ is $c^{a}e_{S}$ for exactly one $a \in \{0, 1\}$ and $S$, so $|G| = 2^{n+1}$, and $G$ is presented by the generators $e_{1}, \dots, e_{n}, c$ and the relations: $c$ central, $c^{2} = 1$, $e_{i}^{2} = c^{s_{i}}$, $e_{i}e_{j} = c\,e_{j}e_{i}$ ($i \ne j$).
\item[(ii)] \emph{(Squares and commutators.)} $e_{S}e_{T} = c^{f(S,T)}e_{S \triangle T}$ with $f(S,T) = |\{(i, j) \in S \times T : i > j\}| + |S \cap T \cap N|$. Every square and every commutator of $G$ lies in $\{1, c\}$; for $g \in \{e_{S}, c\,e_{S}\}$ and $h \in \{e_{T}, c\,e_{T}\}$, $g^{2} = c^{q(S)}$ and $gh = c^{b(S,T)}hg$ with
\[
q(S) = \binom{|S|}{2} + |S \cap N|, \qquad b(S,T) = |S|\,|T| + |S \cap T| \pmod 2,
\]
a quadratic form over $\mathbb{F}_{2}$ and its polar form (the same forms as in Theorem \emph{Blade square signs as a quadratic form over the two-element field}). For $n \ge 2$, $c$ is the only element other than $1$ that is a square or a commutator.
\item[(iii)] \emph{(The centre.)} The radical of $b$ is $\{\emptyset\}$ for $n$ even and $\{\emptyset, \{1, \dots, n\}\}$ for $n$ odd, and the centre of $G$ is $\{1, c\}$ for $n$ even and $\{1, c, e_{1 \cdots n}, c\,e_{1 \cdots n}\}$ for $n$ odd, $e_{1 \cdots n} = e_{1}e_{2}\cdots e_{n}$. So for $n$ even the centre, the commutator subgroup and the subgroup generated by the squares coincide and have order $2$: $G$ is an extraspecial $2$-group of order $2^{1+n}$. The element $c$ is the only central element of order $2$ when $n$ is even or $e_{1 \cdots n}^{2} = c$; when $n$ is odd and $e_{1 \cdots n}^{2} = 1$ (exactly when $p - q \equiv 1 \bmod 4$), $e_{1 \cdots n}$ and $c\,e_{1 \cdots n}$ are central of order $2$ as well.
\item[(iv)] \emph{(The groups are classified by $q$.)} Two groups of signed directions with the same $n$ are isomorphic by an isomorphism carrying $c$ to $c'$ if and only if their forms are equivalent, $q' \circ \phi = q$ for a linear bijection $\phi$ of $\mathbb{F}_{2}^{n}$. For $n \ge 2$ every isomorphism carries $c$ to $c'$.
\item[(v)] \emph{(The integer algebra.)} The ideal $(c + 1)$ is $\bigoplus_{S}\Z\,(1 + c)e_{S}$, and $A$ is a free $\Z$-module with basis the images of the $e_{S}$, written $e_{S}$ again; $c$ maps to $-1$ and $e_{S}e_{T} = (-1)^{f(S,T)}e_{S \triangle T}$, so $A$ has rank $2^{n}$ and integral structure constants.
\item[(vi)] \emph{(The contraction rule as a theorem.)} For $v = \sum_{i} a_{i}e_{i}$ with $a_{i} \in \Z$, $v^{2} = \sum_{i} a_{i}^{2}e_{i}^{2} = Q(v)$.
\item[(vii)] \emph{(The universal property.)} Let $R$ be a commutative ring, $A_{R} = R \otimes_{\Z} A = R[G]/(c + 1)$, free over $R$ with basis the $e_{S}$, and $\iota : R^{n} \to A_{R}$, $x \mapsto \sum_{i} x_{i}e_{i}$. For every associative unital $R$-algebra $B$ and $R$-linear $f : R^{n} \to B$ with $f(x)^{2} = Q(x)\,1$ for all $x$, there is exactly one unital $R$-algebra homomorphism $F : A_{R} \to B$ with $F \circ \iota = f$. So $(A_{R}, \iota)$ is the Clifford algebra $\Cl_{R}(p,q)$ of the quadratic module $(R^{n}, Q)$, isomorphic to the tensor algebra of $R^{n}$ modulo $x \otimes x - Q(x)$, with $\iota$ injective and the blade basis over every $R$; $A = \Cl_{\Z}(p,q)$; and for $R = \R$, $(A_{\R}, \iota)$ is the Clifford algebra $\Cl(p,q)$ of Definition \emph{The Clifford algebra of a quadratic space} [algebra core], the $e_{S}$ its blades. If $2 = 0$ in $R$, $A_{R}$ is the commutative group algebra $R[\mathbb{F}_{2}^{n}]$ of the configurations.
\item[(viii)] \emph{(Isomorphism classes.)} For signatures $(p,q)$ and $(p',q')$ with $p + q = p' + q' = n$ the following are equivalent: (a) $q_{(p,q)}$ and $q_{(p',q')}$ are equivalent over $\mathbb{F}_{2}$; (b) their groups of signed directions are isomorphic by an isomorphism carrying $c$ to $c'$; (c) $\Cl_{R}(p,q) \cong \Cl_{R}(p',q')$ for every commutative ring $R$; (d) $\Cl_{R}(p,q) \cong \Cl_{R}(p',q')$ as $R$-algebras for one commutative ring $R$ that maps unitally to an ordered field ($\Z$, $\Q$ and $\R$ among them); (e) the Gauss sums $\sum_{S}(-1)^{q(S)}$ agree. With $j = e_{1}$ in $\Cl_{\Z}(0,1)$, $\sum_{S}(-1)^{q(S)} = \langle (1 - j)(1 + j)^{p}(1 - j)^{q}\rangle_{0}$, which is $2^{n/2}$, $-2^{n/2}$, $2^{(n+1)/2}$, $-2^{(n+1)/2}$, $0$ for $p - q$ congruent modulo $8$ to $0$ or $2$, $4$ or $6$, $1$, $5$, $3$ or $7$ (the values of Theorem \emph{Blade square signs as a quadratic form over the two-element field}); so (e) holds exactly when $p - q$ and $p' - q'$ lie in the same one of these classes. The Gauss sum is the signature of the trace form $(x, y) \mapsto 2^{-n}\Tr(L_{xy}) = \langle xy\rangle_{0}$ ($L_{z}$ left multiplication), which is diagonal on the blades with entries $e_{S}^{2}$.
\end{enumerate}
\end{theorem}
\begin{proof}
\emph{(i)} For $n \ge 2$ the relation for the pair $(1, 2)$ and the relation for the pair $(2, 1)$ give $e_{1}e_{2} = c\,e_{2}e_{1} = c^{2}e_{1}e_{2}$, so $c^{2} = 1$ by cancellation in $G$ (the same argument as the sign from symmetry, clause (i) of Theorem \emph{Blade square signs as a quadratic form over the two-element field}, read in the group). For $n = 1$ there is no pair. In the cyclic group of order $2m$ generated by $e_{1}$, with $m \ge 2$ and $c = e_{1}^{2}$, the element $c$ is central and $\ne 1$, $G/\langle c\rangle = \{\langle c\rangle, e_{1}\langle c\rangle\}$ has order $2$ with basis the image of $e_{1}$, and $c$ has order $m$, so $c^{2} = 1$ exactly when $m = 2$. Existence: $f(S,T) = \sum_{i > j}x_{i}y_{j} + \sum_{i}s_{i}x_{i}y_{i}$, with $x, y \in \mathbb{F}_{2}^{n}$ the indicator vectors of $S, T$, is $\mathbb{F}_{2}$-bilinear. So on $\{0, 1\} \times \mathbb{F}_{2}^{n}$ the product $(a, S)(b, T) = (a + b + f(S,T), S \triangle T)$ (first entry modulo $2$) is associative, both triple products being $(a + b + d + f(S,T) + f(S,U) + f(T,U), S \triangle T \triangle U)$ for a third factor $(d, U)$; $(0, \emptyset)$ is the identity, as $f$ vanishes when either argument is $\emptyset$; and $(a + f(S,S), S)$ is the inverse of $(a, S)$. Put $c = (1, \emptyset)$, $e_{i} = (0, \{i\})$. Then $c$ is central with $c^{2} = 1 \ne c$; $e_{i}^{2} = (f(\{i\}, \{i\}), \emptyset) = c^{s_{i}}$; for $i \ne j$ the products $e_{i}e_{j}$ and $e_{j}e_{i}$ are $(f(\{i\},\{j\}), \{i, j\})$ and $(f(\{j\},\{i\}), \{i, j\})$, whose first entries add to $1$, so $e_{i}e_{j} = c\,e_{j}e_{i}$; and since $f(S, \{j\}) = 0$ when $j$ exceeds every element of $S$, the increasing product $e_{S}$ is $(0, S)$, so $(a, S) \mapsto S$ identifies the quotient by $\langle c\rangle$ with $\mathbb{F}_{2}^{n}$, the $e_{i}$ mapping to the basis. This is a group of signed directions with the given signs; uniqueness up to an isomorphism fixing $c$ and the $e_{i}$ follows from the presentation below. Granting $c^{2} = 1$ and $c \ne 1$, $\langle c\rangle = \{1, c\}$; each coset of $\langle c\rangle$ is $e_{S}\langle c\rangle$ for exactly one $S$, so each element is $c^{a}e_{S}$ for exactly one $a \in \{0, 1\}$ and one $S$, and $|G| = 2^{n+1}$. The image of $e_{i}^{2}$ in $\mathbb{F}_{2}^{n}$ is $\{i\} \triangle \{i\} = \emptyset$, so $e_{i}^{2} \in \{1, c\}$. Let $\Gamma$ be the group presented by the generators and relations of the statement. In $\Gamma$, $c^{-1} = c$ and $e_{i}^{-1} = c^{s_{i}}e_{i}$, so every element is a word in $c, e_{1}, \dots, e_{n}$ without inverses, and the rewriting of (ii), which uses only the relations, brings it to the form $c^{a}e_{S}$; hence $|\Gamma| \le 2^{n+1}$. The elements $e_{i}$ and $c$ of $G$ satisfy the relations and generate $G$, so the homomorphism $\Gamma \to G$ sending generators to generators is onto, and since $|G| = 2^{n+1} \ge |\Gamma|$ it is an isomorphism.

\emph{(ii)} In the word $e_{S}e_{T} = e_{i_{1}}\cdots e_{i_{k}}e_{j_{1}}\cdots e_{j_{l}}$ ($i_{1} < \dots < i_{k}$, $j_{1} < \dots < j_{l}$) move $e_{j_{1}}, e_{j_{2}}, \dots$ in turn to the left until they reach their place in increasing order. The letter $e_{j}$ passes exactly the letters $e_{i}$ of $e_{S}$ with $i > j$, each passage costing a factor $c$ (the relation, with $c$ central), and if $j \in S$ it then meets $e_{j}$ and the two merge into $e_{j}^{2} = c^{s_{j}}$; it never meets a letter of $e_{T}$ already placed, those being smaller. Collecting the central factors gives $e_{S}e_{T} = c^{f(S,T)}e_{S \triangle T}$, the sign rule of the blade product of an orthogonal basis, with $c$ in place of $-1$. Since $G/\langle c\rangle$ is elementary abelian, every square $g^{2}$ and every commutator $ghg^{-1}h^{-1}$ maps to $0$ and lies in $\langle c\rangle = \{1, c\}$, which is central. For $g = c^{a}e_{S}$, $g^{2} = c^{2a}e_{S}^{2} = c^{f(S,S)}$, and $f(S,S) = \binom{|S|}{2} + |S \cap N|$, the pairs $(i, j) \in S \times S$ with $i > j$ being the $\binom{|S|}{2}$ two-element subsets of $S$. Also $e_{T}e_{S} = c^{f(T,S)}e_{S \triangle T}$, so, as $c^{-1} = c$, $e_{S}e_{T} = c^{f(S,T) + f(T,S)}e_{T}e_{S}$, and $f(S,T) + f(T,S)$ counts the ordered pairs $(i, j) \in S \times T$ with $i \ne j$, plus $2|S \cap T \cap N|$: it is $|S|\,|T| - |S \cap T| \equiv |S|\,|T| + |S \cap T| = b(S,T) \bmod 2$. The central factors $c^{a}$ do not change squares or commutators. From $hg = c^{b(S,T)}gh$, $(gh)^{2} = g(hg)h = c^{b(S,T)}g^{2}h^{2}$, so $q(S \triangle T) = q(S) + q(T) + b(S,T)$: $q$ is a quadratic form over $\mathbb{F}_{2}$ with polar form $b = f + f^{\top}$, which is bilinear because $f$ is; these are the same forms as in clause (ii) of Theorem \emph{Blade square signs as a quadratic form over the two-element field} (there with the negative directions numbered last; the formulas do not depend on the numbering). For $n \ge 2$, $e_{1}e_{2}e_{1}^{-1}e_{2}^{-1} = c$ is a commutator. It is also a square: if $s_{i} = s_{j}$ for some $i \ne j$, then $(e_{i}e_{j})^{2} = c^{q(\{i, j\})} = c^{1 + s_{i} + s_{j}} = c$; otherwise $n = 2$ and $\{s_{1}, s_{2}\} = \{0, 1\}$, and $e_{i}^{2} = c$ for the $i$ with $s_{i} = 1$.

\emph{(iii)} An element $c^{a}e_{S}$ is central exactly when it commutes with every $e_{i}$, since the $e_{i}$ and $c$ generate $G$, that is, when $b(S, \{i\}) = 0$ for every $i$: when $S$ lies in the radical of $b$. Here $b(S, \{i\}) = |S| + [i \in S]$. If $S$ is neither empty nor $\{1, \dots, n\}$, taking $i \in S$ and $i' \notin S$ gives $|S| \equiv 1$ and $|S| \equiv 0$, which is impossible; $\emptyset$ is always in the radical, and $\{1, \dots, n\}$ is exactly when $n + 1 \equiv 0$, that is $n$ odd. So the radical is $\{\emptyset\}$ for $n$ even and $\{\emptyset, \{1, \dots, n\}\}$ for $n$ odd (as in clause (iii) of Theorem \emph{Blade square signs as a quadratic form over the two-element field}), which gives the centre. For $n$ even, the centre, the commutators and the squares are all $\{1, c\}$ by (ii) ($n \ge 2$), so the subgroup they generate is $\{1, c\}$ as well, and $G$ is extraspecial of order $2^{1+n}$ (for a finite $2$-group the subgroup generated by squares and commutators is the Frattini subgroup). Central elements of order $2$: $c$, and for $n$ odd $c^{a}e_{1 \cdots n}$ exactly when $(c^{a}e_{1 \cdots n})^{2} = c^{q(\{1, \dots, n\})} = 1$. Here $q(\{1, \dots, n\}) = \binom{n}{2} + q$, and for $n = p + q$ odd, $\binom{n}{2} + q$ is even exactly when $p - q \equiv 1 \bmod 4$: writing $p - q = n - 2q$, both conditions are unchanged when $q$ grows by $2$ and both flip when $q$ grows by $1$ (since $n - 2q$ then drops by $2$), and for $q = 0$ the first reads $n \equiv 1 \bmod 4$, which is the second.

\emph{(iv)} Let $\Phi : G \to G'$ be an isomorphism with $\Phi(c) = c'$. It induces a group isomorphism $\phi : G/\langle c\rangle \to G'/\langle c'\rangle$, that is, a linear bijection of $\mathbb{F}_{2}^{n}$, and $\Phi(g)^{2} = \Phi(g^{2})$ reads $c'^{\,q'(\phi(S))} = c'^{\,q(S)}$, so $q' \circ \phi = q$. Conversely, let $q' \circ \phi = q$. Then also $b' \circ (\phi \times \phi) = b$, $b$ being the polar form of $q$. Choose $g_{i} \in G'$ over $\phi(\{i\})$. By (ii) for $G'$, $g_{i}^{2} = c'^{\,q'(\phi(\{i\}))} = c'^{\,s_{i}}$ and $g_{i}g_{j} = c'^{\,b(\{i\}, \{j\})}g_{j}g_{i} = c'\,g_{j}g_{i}$ for $i \ne j$, and $c'$ is central with $c'^{2} = 1$. By the presentation of (i) there is a homomorphism $\Phi : G \to G'$ with $\Phi(e_{i}) = g_{i}$ and $\Phi(c) = c'$; its image contains $c'$ and elements over the basis $\phi(\{1\}), \dots, \phi(\{n\})$, hence every coset of $\langle c'\rangle$ and all of $G'$; so $\Phi$ is onto, and an isomorphism because $|G| = |G'| = 2^{n+1}$. For $n \ge 2$, $c$ is the only commutator other than $1$ in $G$ and $c'$ in $G'$, and isomorphisms carry commutators to commutators.

\emph{(v)} $\Z[G]$ is free on the $2^{n+1}$ elements $c^{a}e_{S}$. As $c$ is central, the two-sided ideal generated by $c + 1$ is $(c + 1)\Z[G]$, which is spanned over $\Z$ by the $(c + 1)c^{a}e_{S}$; since $c(c + 1) = c + c^{2} = c + 1$, these are the $(1 + c)e_{S}$. So $\Z[G] = \bigoplus_{S}(\Z e_{S} \oplus \Z\,c\,e_{S})$ and $(c + 1) = \bigoplus_{S}\Z(e_{S} + c\,e_{S})$; in each summand $e_{S}$ and $e_{S} + c\,e_{S}$ form a $\Z$-basis (the change of basis from $e_{S}, c\,e_{S}$ has determinant $1$), so $A = \bigoplus_{S}\Z e_{S}$, free with basis the images of the $e_{S}$, with $c\,e_{S} \mapsto -e_{S}$. The product follows from (ii).

\emph{(vi)} $v^{2} = \sum_{i}a_{i}^{2}e_{i}^{2} + \sum_{i < j}a_{i}a_{j}(e_{i}e_{j} + e_{j}e_{i})$, and $e_{i}e_{j} + e_{j}e_{i} = (c + 1)e_{j}e_{i}$ maps to $0$ in $A$, while $e_{i}^{2} = c^{s_{i}}$ maps to $(-1)^{s_{i}}$. So $v^{2} = \sum_{i}(-1)^{s_{i}}a_{i}^{2} = Q(v)$; no element $\tfrac12$ is used.

\emph{(vii)} Tensoring the exact sequence $(c + 1) \to \Z[G] \to A \to 0$ with $R$ gives $A_{R} = R[G]/(c + 1)$, free over $R$ on the $1 \otimes e_{S}$, which are products of the $\iota(\epsilon_{i}) = e_{i}$ ($\epsilon_{i}$ the standard basis of $R^{n}$). Uniqueness: $A_{R}$ is generated as an $R$-algebra by $\iota(R^{n})$. Existence: put $u_{i} = f(\epsilon_{i})$. Then $u_{i}^{2} = Q(\epsilon_{i})\,1 = (-1)^{s_{i}}$, and for $i \ne j$, $(u_{i} + u_{j})^{2} = Q(\epsilon_{i} + \epsilon_{j}) = (-1)^{s_{i}} + (-1)^{s_{j}}$ gives $u_{i}u_{j} + u_{j}u_{i} = 0$. Each $u_{i}$ is a unit, with inverse $(-1)^{s_{i}}u_{i}$, and in the unit group of $B$ the elements $u_{1}, \dots, u_{n}$ and $-1$ satisfy the relations of the presentation in (i): $-1$ is central with square $1$, $u_{i}^{2} = (-1)^{s_{i}}$, $u_{i}u_{j} = (-1)u_{j}u_{i}$. So there is a group homomorphism $\psi : G \to B^{\times}$ with $\psi(e_{i}) = u_{i}$ and $\psi(c) = -1$; its $R$-linear extension $R[G] \to B$ is a unital algebra homomorphism that kills $c + 1$, hence factors through $A_{R}$ as $F$, and $F(\iota(x)) = \sum_{i}x_{i}u_{i} = f(x)$. The tensor algebra of $R^{n}$ modulo the ideal generated by the $x \otimes x - Q(x)$ has the same universal property (that of the tensor algebra, passed to the quotient), so the two are isomorphic by the unique isomorphism compatible with the maps from $R^{n}$; $\iota$ is injective because the $e_{i}$ are part of an $R$-basis. For $R = \R$ the universal property is that of Definition \emph{The Clifford algebra of a quadratic space} [algebra core] for $\R^{n}$ with the form $Q$ of signature $(p, q)$. If $2 = 0$ in $R$ then $-1 = 1$, so $e_{S}e_{T} = e_{S \triangle T}$, which is the multiplication of the group algebra of $\mathbb{F}_{2}^{n}$.

\emph{(viii)} (a) $\Leftrightarrow$ (b) is (iv). (b) $\Rightarrow$ (c): an isomorphism $G \to G'$ carrying $c$ to $c'$ extends to $R[G] \cong R[G']$ carrying $c + 1$ to $c' + 1$. (c) $\Rightarrow$ (d) with $R = \Z$. (d) $\Rightarrow$ (e): let $R \to K$ be a unital homomorphism to an ordered field; then $K \otimes_{R} A_{R} = A_{K}$, so an isomorphism over $R$ gives one over $K$. On $A_{K}$ put $\tau(x, y) = 2^{-n}\Tr(L_{xy})$, the trace of left multiplication on the $K$-vector space $A_{K}$. An algebra isomorphism $\Phi$ satisfies $L_{\Phi(z)} = \Phi L_{z}\Phi^{-1}$, so $\tau$ is carried to $\tau'$. Since $e_{S}e_{T} = \pm e_{S \triangle T}$, $L_{e_{S}}$ maps each basis line to a different one unless $S = \emptyset$, so $\Tr(L_{e_{S}}) = 2^{n}$ for $S = \emptyset$ and $0$ otherwise; hence $\tau(x, y) = \langle xy\rangle_{0}$, $\tau(e_{S}, e_{T}) = 0$ for $S \ne T$, and $\tau(e_{S}, e_{S}) = (-1)^{q(S)}$. Let $P$ be the number of $S$ with $q(S) = 0$ and $U_{\pm}$ the spans of the blades with $e_{S}^{2} = \pm 1$. On $U_{+}$, $\tau(x, x) = \sum x_{S}^{2} > 0$ for $x \ne 0$, and on $U_{-}$, $\tau(x, x) = -\sum x_{S}^{2} < 0$ for $x \ne 0$; a subspace $W$ with $\dim W > P$ meets $U_{-}$ in a nonzero vector, as $\dim W + \dim U_{-} > 2^{n}$. So $P$ is the largest dimension of a subspace on which $\tau$ is positive definite, an isomorphism invariant, and with it the Gauss sum $\sum_{S}(-1)^{q(S)} = 2P - 2^{n}$. (e) $\Rightarrow$ (a): all the forms $q_{(p,q)}$ with the same $n$ have the same polar form $b$, whose radical is $\{0\}$ for $n$ even and $\{0, z\}$, $z = \{1, \dots, n\}$, for $n$ odd. Let $W$ be a complement of the radical ($W = \mathbb{F}_{2}^{n}$ for $n$ even); $b$ is alternating ($b(x, x) = q(0) = 0$) and nondegenerate on $W$, so $W$ has a basis $u_{1}, w_{1}, \dots, u_{m}, w_{m}$ with $b(u_{k}, w_{k}) = 1$ and $b = 0$ on all other pairs of basis vectors (choose $u \ne 0$, then $w$ with $b(u, w) = 1$, and continue in the $b$-orthogonal complement of $\mathrm{span}\{u, w\}$, which is complementary to it and on which $b$ is nondegenerate). On each plane $\mathrm{span}\{u_{k}, w_{k}\}$ the three nonzero values of $q$ add to $1$, since $q(u + w) = q(u) + q(w) + 1$: either two of them are $0$, and the two vectors with $q = 0$ form a new basis of the plane with $b = 1$ (type $H$, Gauss sum $2$), or all three are $1$ (type $P$, Gauss sum $-2$). Gauss sums multiply over $b$-orthogonal sums, since $q(x + y) = q(x) + q(y)$ when $b(x, y) = 0$. Two planes of type $P$ give one of type $H$ and one more: in $\mathrm{span}\{u_{1}, w_{1}, u_{2}, w_{2}\}$, $u = u_{1} + u_{2}$ and $w = w_{1} + u_{2}$ have $q(u) = q(w) = 0$ and $b(u, w) = 1$, and the orthogonal complement of $\mathrm{span}\{u, w\}$ there is a plane of Gauss sum $(-2)(-2)/2 = 2$, so of type $H$. Hence $q|_{W}$ is equivalent to $H^{m}$, Gauss sum $2^{m}$, or to $H^{m-1} \perp P$, Gauss sum $-2^{m}$. For $n$ even this is $q$ itself. For $n$ odd, $q(x + az) = q(x) + a\,q(z)$ for $x \in W$, as $z$ is in the radical. If $q(z) = 1$, the pairing $x \leftrightarrow x + z$ shows that the Gauss sum is $0$, and a plane of type $P$ in $W$ becomes one of type $H$ when $u$ is replaced by $u + z$ ($q(u + z) = 0$, $b$ unchanged), so $q$ is equivalent to $H^{m} \perp \langle z\rangle$ with $q(z) = 1$; if $q(z) = 0$, $q$ is equivalent to $q|_{W} \perp \langle z\rangle$ with $q(z) = 0$, with Gauss sum $\pm 2^{m+1}$, the sign that of $q|_{W}$. In each parity the models have distinct Gauss sums, so two forms with equal Gauss sums are equivalent to the same model, hence to each other. The values: in $\Cl_{\Z}(0,1) = \Z \oplus \Z j$, $j^{2} = -1$, a commutative ring, $(-1)^{\binom{k}{2}} = \langle (1 - j)j^{k}\rangle_{0}$ (both sides are $1, 1, -1, -1$ for $k = 0, 1, 2, 3$ and have period $4$). Grouping the $S$ by $a = |S \setminus N|$ and $b' = |S \cap N|$, $(-1)^{q(S)} = (-1)^{\binom{a + b'}{2} + b'} = \langle (1 - j)j^{a}(-j)^{b'}\rangle_{0}$, so $\sum_{S}(-1)^{q(S)} = \langle (1 - j)(1 + j)^{p}(1 - j)^{q}\rangle_{0}$. With $(1 + j)(1 - j) = 2$, $(1 + j)^{2} = 2j$ and $(1 + j)^{8} = 16$: for $r = p - q = 8m + s \ge 0$ ($0 \le s \le 7$) the product is $2^{q}16^{m}(1 - j)(1 + j)^{s}$, and $\langle (1 - j)(1 + j)^{s}\rangle_{0} = 1, 2, 2, 0, -4, -8, -8, 0$ for $s = 0, \dots, 7$; for $r < 0$, $|r| = 8m + s'$, the product is $2^{p}(1 - j)^{|r| + 1}$, whose scalar part equals that of $2^{p}(1 + j)^{|r| + 1}$ (the automorphism $j \mapsto -j$), and $\langle (1 + j)^{t}\rangle_{0} = 1, 1, 0, -2, -4, -4, 0, 8, 16$ for $t = 0, \dots, 8$. With $n = r + 2q$ (respectively $|r| + 2p$) this is $2^{n/2}$ for $r \equiv 0, 2$, $-2^{n/2}$ for $r \equiv 4, 6$, $2^{(n+1)/2}$ for $r \equiv 1$, $-2^{(n+1)/2}$ for $r \equiv 5$ and $0$ for $r \equiv 3, 7 \pmod 8$, the values of clause (iv) of Theorem \emph{Blade square signs as a quadratic form over the two-element field}; for fixed $n$ they are distinct on the five classes, which are therefore the classes of (viii).

The check builds the cocycle group for every sign vector with $n \le 4$ and verifies associativity, inverses and the relations; it builds $G$ from its relations by rewriting words for every sign vector with $n \le 4$ and for the standard signatures with $n = 5$, and verifies the normal form, associativity ($n \le 3$), agreement with the product of signed blades in GA with $c \mapsto -1$, the closed forms of $q$ and $b$, and that the squares are $\{1, c\}$; the Heisenberg group modulo $k$ ($2 \le k \le 7$) reads its relation both ways only for $k = 2$, and the cyclic groups of order $2m$ ($2 \le m \le 7$) show the independence at $n = 1$; the centre, the central involutions and the extraspecial property for all $20$ signatures with $n \le 5$; the ideal, a unimodular basis change, integral structure constants equal to GA's and $v^{2} = Q(v)$ for all integer vectors with entries in $[-2, 2]$, $n \le 4$, and rational vectors; the universal property for $27$ maps into $\Cl(p{+}1, q{+}1)$ (twisted generators), rotor-conjugated frames and the grade involution; the reduction modulo $2$; and, for all pairs of sign vectors with $n \le 3$ and of standard signatures with $4 \le n \le 6$, that an explicit $\mathbb{F}_{2}$-isometry exists exactly when the Gauss sums agree, lifted for $n \le 4$ to anticommuting blades of the right squares whose products are $2^{n}$ distinct signed blades (an isomorphism over $\Z$), with the traces $\Tr(L_{e_{S}})$ and the diagonal of the trace form; and the radical and the scalar part of $(1 - j)(1 + j)^{p}(1 - j)^{q}$ in $\Cl_{\Z}(0,1)$, equal to the Gauss sum and to the tabulated value, for all signatures with $n \le 8$.

\end{proof}

\phantomsection\label{c:note.geometry-first.setting-superposition-one-rule-theorem-3}
\noindent\emph{How the proof goes.} Setting $c = -1$ leaves the $2^{n}$ blades as a basis, and in $v^{2}$ the cross terms cancel in pairs because different directions anticommute, which leaves $Q$. Part (iii) follows by running the same sorting in any algebra whose vectors square by $Q$. Example~\ref{c:note.geometry-first.one-rule-hand} is the case $n = 2$. The proof is in Appendix~\ref{prf:thm:bitclifford}.

The primer \cite[Section 2]{ButlerT0} takes $v^{2} = Q(v)$ as the one rule from which the geometric product is generated. From the bit it is not a rule but a consequence: it is what one bit of non-commutativity and the freedom to add leave behind. The construction of a Clifford algebra as a group algebra twisted by a sign is known \cite{AlbuquerqueMajid2002}; what is added here is its reading as the second and third steps from a single distinction.

\begin{example}[The one rule by hand]\label{c:note.geometry-first.one-rule-hand}
Take $n = 2$ with $N = \{2\}$, so that after (F3) $e_{1}^{2} = 1$, $e_{2}^{2} = -1$ and $e_{1}e_{2} = -e_{2}e_{1}$. For integers $a$, $b$ and $v = a e_{1} + b e_{2}$,
\[
v^{2} \;=\; a^{2}e_{1}^{2} + ab\,(e_{1}e_{2} + e_{2}e_{1}) + b^{2}e_{2}^{2} \;=\; a^{2} - b^{2} \;=\; Q(v) .
\]
The cross terms cancel because the two directions anticommute, and the minus sign is the choice $N = \{2\}$. Nothing about $v^{2}$ was assumed: the one rule came out. With $a = b = 1$ the vector $e_{1} + e_{2}$ squares to $0$, a null vector over the integers.
\end{example}

\section{Division, and the rationals}

\phantomsection\label{c:note.geometry-first.setting-division-rationals}
The next step is a demand rather than a premise about the bit. It is the primer's requirement \cite[Section 2]{ButlerT0} that the product can be undone, asked of the integer vectors that can be inverted at all:

\begin{proposition}[Division in the integer Clifford algebra gives the rationals]\label{c:synthesis.division-gives-the-rationals}
\changedmark{BurntOrange}{changed since the deposit}
Let $n = p + q \ge 1$, $N = \{p + 1, \dots, n\}$, and for a commutative ring $R$ let $\Cl_{R}(p,q)$ be the free $R$-module with basis the blades $e_{S}$, $S \subseteq \{1, \dots, n\}$ ($e_{\emptyset} = 1$), with the blade product
\[
e_{S}e_{T} = (-1)^{f(S,T)}\,e_{S \triangle T}, \qquad f(S,T) = \#\{(i,j) \in S \times T : i > j\} + |S \cap T \cap N| ,
\]
the same rule as Proposition \ref{c:std-algebra.orthogonal-anticommute} gives over the reals for squares $Q(e_{i}) = 1$ ($i \le p$), $Q(e_{i}) = -1$ ($i \in N$). It is an associative unital $R$-algebra, $\Cl_{R}(p,q) = R \otimes_{\Z} \Cl_{\Z}(p,q)$, and for $v = \sum_{i} a_{i}e_{i}$ in $V_{R} = \bigoplus_{i} R\,e_{i}$, $v^{2} = Q(v) = \sum_{i \le p}a_{i}^{2} - \sum_{i \in N}a_{i}^{2}$.
\begin{enumerate}
\item[(i)] \emph{(Inverses.)} For a commutative ring $R$ and $v \in V_{R}$: $v$ has a left or a right inverse in $\Cl_{R}(p,q)$ if and only if $Q(v)$ is a unit of $R$, and then $v^{-1} = Q(v)^{-1}v$ is a two-sided inverse (the formula of Proposition \emph{Which vectors and versors are invertible} [algebra core]).
\item[(ii)] \emph{(Division produces $\Q$.)} For a commutative ring $R$, every $v \in V_{\Z}$ with $Q(v) \ne 0$ becomes invertible in $\Cl_{R}(p,q)$ if and only if every nonzero integer is a unit of $R$, that is, $R$ is a $\Q$-algebra. So $\Q$ is the initial such ring, $\Cl_{\Q}(p,q)$ is the localisation of $\Cl_{\Z}(p,q)$ at the nonzero integers, and no subring of $\Q$ other than $\Q$ itself will do. Already the multiples $m e_{1}$, $m \ge 1$, force it: $Q(me_{1}) = \pm m^{2}$ and $e_{1}(me_{1})^{-1} = 1/m$, so every rational $a/m$ is the scalar $(ae_{1})(me_{1})^{-1}$, a product of an integer vector and the inverse of one. In $\Cl_{\Q}(p,q)$ every $v \in V_{\Q}$ with $Q(v) \ne 0$ is invertible, and the grade-$0$ part is $\Q\,1$.
\item[(iii)] \emph{(What division leaves.)} $V_{\Q}$ contains a nonzero null vector ($v^{2} = 0$, a zero divisor of $\Cl_{\Q}(p,q)$) if and only if $p \ge 1$ and $q \ge 1$, for example $e_{1} + e_{p+1}$. The algebra $\Cl_{\Q}(p,q)$ contains zero divisors if and only if $(p, q) \notin \{(0,0), (0,1), (0,2)\}$, with integral witnesses $(1 + e_{1})(1 - e_{1}) = 0$ for $p \ge 1$ and $(1 + e_{1}e_{2}e_{3})(1 - e_{1}e_{2}e_{3}) = 0$ for $p = 0$, $q \ge 3$; for $(p, q) \in \{(0,0), (0,1), (0,2)\}$, $x\bar{x} = \bar{x}x$ is the sum of the squares of the coefficients of $x$ ($\bar{x}$ the Clifford conjugate), so $\Cl_{\Q}(p,q)$ is a division ring. This is the list of Theorem \ref{c:transport.exceptional-division-algebras}, now over $\Q$.
\end{enumerate}
\end{proposition}
\begin{proof}
\emph{The ring.} The function $f$ is bilinear over $\mathbb{F}_{2}$ in the indicator vectors of $S$ and $T$ (each term counts pairs, and a pair counted in $S \triangle S'$ is counted in exactly one of $S$, $S'$ modulo $2$). Hence $f(S \triangle T, U) + f(S, T) \equiv f(S, U) + f(T, U) + f(S, T) \equiv f(S, T \triangle U) + f(T, U) \pmod 2$, which says $(e_{S}e_{T})e_{U} = e_{S}(e_{T}e_{U})$; extended bilinearly the product is associative, with unit $e_{\emptyset}$, and the structure constants are $\pm 1$, so $\Cl_{\Z}(p,q)$ is a ring and, the same computation holding over any commutative ring $R$, $\Cl_{R}(p,q) = R \otimes_{\Z}\Cl_{\Z}(p,q)$ is an $R$-algebra with $R$ central. The rule is the one Proposition \ref{c:std-algebra.orthogonal-anticommute} gives over the reals, with $\prod_{i \in S \cap T}Q(e_{i}) = (-1)^{|S \cap T \cap N|}$; it is used here as a comparison, the ring being defined by the rule. For $v = \sum a_{i}e_{i}$, $v^{2} = \sum_{i}a_{i}^{2}e_{i}^{2} + \sum_{i < j}a_{i}a_{j}(e_{i}e_{j} + e_{j}e_{i}) = Q(v)$, since $e_{j}e_{i} = -e_{i}e_{j}$ for $i \ne j$ ($f(\{j\}, \{i\}) - f(\{i\}, \{j\}) = \pm 1$) and $e_{i}^{2} = (-1)^{[i \in N]}$. The check verifies the product rule against GA.py for every pair of blades in all signatures with $n \le 4$, the bilinearity on all triples, and $v^{2} = Q(v)$ on a box of integer vectors.

\emph{(i)} If $Q(v)$ is a unit of $R$, put $u = Q(v)^{-1}v \in \Cl_{R}(p,q)$; then $vu = uv = Q(v)^{-1}v^{2} = 1$, the formula of Proposition \emph{Which vectors and versors are invertible} [algebra core], here proved over $R$. Conversely, if $vw = 1$ for some $w$, then $v^{2}w^{2} = v(vw)w = vw = 1$, that is $Q(v)\,w^{2} = 1$ since $v^{2} = Q(v)1$ is central; if $wv = 1$, likewise $w^{2}v^{2} = 1$. Comparing grade-$0$ coefficients in the free $R$-module, $Q(v)\,\langle w^{2}\rangle_{0} = 1$, so $Q(v)$ is a unit of $R$, and the one-sided inverse equals $u$ by associativity ($w = (uv)w = u(vw) = u$, or $w = w(vu) = (wv)u = u$). The check verifies, for all integer vectors with coordinates in $\{-3, \dots, 3\}$ and $Q(v) \ne 0$ in $\Cl(2,0)$, $\Cl(1,1)$, $\Cl(0,3)$, $\Cl(1,2)$, that $v/Q(v)$ is a two-sided inverse and that its coefficients lie in $\Z[1/P]$ exactly when $Q(v)$ is a unit there, for $P = \emptyset, \{2\}, \{3\}, \{2,3\}, \{2,5\}$, counts the vectors invertible over $\Z$ against a hand count ($Q(v) = \pm 1$), and, over the ring $\Z/6\Z$ (not a subring of $\Q$), finds by exhaustive search in $\Cl(1,1)$ that a vector has an inverse exactly when $Q(v)$ is prime to $6$.

\emph{(ii)} If every nonzero integer is a unit of $R$, every $v \in V_{\Z}$ with $Q(v) \ne 0$ is invertible by (i). Conversely, if every such $v$ is invertible in $\Cl_{R}(p,q)$, take $v = me_{1}$, $m \ge 1$: $Q(v) = m^{2}Q(e_{1}) = \pm m^{2} \ne 0$, so by (i) $m^{2}$ is a unit of $R$, and so is $m$ ($m \cdot m\,(m^{2})^{-1} = 1$). A commutative ring in which every nonzero integer is a unit receives exactly one ring homomorphism from $\Q$, $a/m \mapsto a\,m^{-1}$, so it is a $\Q$-algebra and $\Q$ is initial among such rings; $\Cl_{\Q}(p,q) = \Q \otimes_{\Z} \Cl_{\Z}(p,q)$ is the localisation of $\Cl_{\Z}(p,q)$ at the nonzero integers, which are central. A subring of $\Q$ that is a $\Q$-algebra contains $1/m$ for all $m$, so it is $\Q$. Explicitly $(me_{1})^{-1} = me_{1}/(m^{2}Q(e_{1})) = e_{1}/(mQ(e_{1}))$ and $e_{1}(me_{1})^{-1} = e_{1}^{2}/(mQ(e_{1})) = 1/m$, so $a/m = (ae_{1})(me_{1})^{-1}$. A vector of $V_{\Q}$ with $Q(v) \ne 0$ is invertible over $\Q$ by (i). The grade-$0$ part of the free $\Q$-module $\Cl_{\Q}(p,q)$ is $\Q e_{\emptyset}$. The check computes $e_{1}(me_{1})^{-1} = 1/m$ and $(ae_{1})(me_{1})^{-1} = a/m$ for $m = 1, \dots, 30$ in three signatures, and the primes that must be inverted, $2, 3, 5, \dots, 29$.

\emph{(iii)} If $p, q \ge 1$, $(e_{1} + e_{p+1})^{2} = 1 - 1 = 0$. If $p = 0$ or $q = 0$, $Q$ is plus or minus a sum of squares on $V_{\Q}$ and vanishes only at $v = 0$. A nonzero $v$ with $v^{2} = 0$ is a zero divisor. If $p \ge 1$, $(1 + e_{1})(1 - e_{1}) = 1 - e_{1}^{2} = 0$ with both factors nonzero. If $p = 0$ and $q \ge 3$, $w = e_{1}e_{2}e_{3}$ has $w^{2} = (-1)^{3}e_{1}^{2}e_{2}^{2}e_{3}^{2} = (-1)(-1)^{3} = 1$ (reversing three letters costs three transpositions), so $(1 + w)(1 - w) = 0$. The remaining signatures are $(0,0)$, $(0,1)$, $(0,2)$. The Clifford conjugate $\bar{x}$ (grade involution after reversion) negates every blade of grade $1$ and $2$; for $x = \sum_{S}x_{S}e_{S}$ in these algebras $e_{S}\bar{e}_{S} = 1$ for every blade ($e_{1}(-e_{1}) = 1$, $e_{1}e_{2}(-e_{1}e_{2}) = 1$), and the cross terms cancel in pairs: for $S \ne T$, $x_{S}x_{T}(e_{S}\bar{e}_{T} + e_{T}\bar{e}_{S}) = 0$, because $\bar{e}_{T} = -e_{T}$ when $T \ne \emptyset$ and $e_{S}e_{T} = -e_{T}e_{S}$ for distinct nonempty $S, T$ in $\Cl(0,2)$, while for $S = \emptyset$ the pair is $x_{\emptyset}x_{T}(\bar{e}_{T} + e_{T}) = 0$. So $x\bar{x} = \sum_{S}x_{S}^{2}$, likewise $\bar{x}x$, which is a positive rational for $x \ne 0$; then $x^{-1} = \bar{x}/(x\bar{x})$ and $\Cl_{\Q}(p,q)$ is a division ring. The list is that of Theorem \ref{c:transport.exceptional-division-algebras}. The check confirms, for every signature with $n \le 5$, the existence of nonzero null vectors exactly for $p, q \ge 1$ (on a box of coordinates), the integral zero divisors, and $x\bar{x} = \bar{x}x = \sum_{S}x_{S}^{2}$ on a grid of rational elements of $\Cl(0,0)$, $\Cl(0,1)$, $\Cl(0,2)$, the only signatures without a witness.

\end{proof}

\phantomsection\label{c:note.geometry-first.setting-division-rationals-2}
\noindent\emph{How the proof goes.} A vector with $Q(v) \neq 0$ has the inverse $v/Q(v)$; conversely the vector $m e_{1}$ shows that each nonzero integer $m$ must be invertible. Example~\ref{c:note.geometry-first.vector-needs-half} is the smallest case. The proof is in Appendix~\ref{prf:prop:bitdivision}.

Every nonzero ring in which every nonzero integer is invertible contains exactly one copy of $\Q$. So the rational numbers are what the demand (D) forces on the coefficients, and nothing less will do: every fraction $a/m$ is the product $(a e_{1})(m e_{1})^{-1}$ of an integer vector and the inverse of one. Premise (F4) below posits a field containing $\Q$ in any case; what (D) and Proposition~\ref{c:synthesis.division-gives-the-rationals} add is that nothing smaller than $\Q$ will do once division is asked for.

Division does not remove everything. When $p \geq 1$ and $q \geq 1$ the null vectors of the primer \cite[Section 4]{ButlerT0} are there over $\Z$ and remain over $\Q$, for instance $e_{i} + e_{j}$ with $i \notin N$ and $j \in N$, and no extension of the coefficients inverts them. That is where T1 begins.

\begin{example}[A vector that needs a half]\label{c:note.geometry-first.vector-needs-half}
With $n = 2$ and $N$ empty, the integer vector $v = e_{1} + e_{2}$ has $v^{2} = 2$, so
\[
v \cdot \tfrac12\,v \;=\; 1 ,
\]
and its inverse $\tfrac12(e_{1} + e_{2})$ exists over $\Q$ but not over $\Z$. In $\Cl(1,1)$, by contrast, the null vector $e_{1} + e_{2}$ of Example~\ref{c:note.geometry-first.one-rule-hand} has no inverse over any nonzero coefficient ring: if $vu = 1$ then $v = v(vu) = v^{2}u = 0$, which is false.
\end{example}

\section{Completion, and the number systems inside}

\phantomsection\label{c:note.geometry-first.setting-completion-number-systems-inside}
This is the one step that is not algebraic, and it is the last premise. What it gives is proved here in three stages: square roots, and with them the classification of quadratic forms by their signs; roots of polynomials; and the list of the algebras in which division never fails. The results are classical \cite{Lam2005, Ebbinghaus1991}; the proofs are written out so that the step to the real numbers rests on (F4) and on nothing else.

Some vocabulary first. Let $F$ be an ordered field. A \emph{quadratic space} over $F$ is a finite-dimensional vector space $V$ with a map $Q : V \to F$ such that $Q(\lambda v) = \lambda^{2}Q(v)$ and
\[
B(u, v) \;=\; \tfrac12\bigl(Q(u + v) - Q(u) - Q(v)\bigr)
\]
is bilinear; then $Q(v) = B(v, v)$. Vectors $u, v$ are \emph{orthogonal} when $B(u,v) = 0$; in any algebra in which $v^{2} = Q(v)$, expanding $(u+v)^{2}$ gives $uv + vu = 2B(u,v)$, so orthogonal vectors are exactly those that anticommute. An \emph{isometry} is a linear bijection that preserves $Q$. The field $F$ is \emph{Euclidean} when every positive element is a square (and $0 = 0^{2}$, so then every element that is not negative is one). For $a \in F$, $|a|$ is $a$ or $-a$, whichever is not negative; one checks from the rules of the order, case by case on the signs, that $|ab| = |a|\,|b|$ and $|a + b| \leq |a| + |b|$. Finally, in an ordered field $1 = 1^{2} > 0$, so $m \cdot 1 > 0$ for every positive integer $m$: the field contains $\Q$, it is infinite, and division by $2$ is always possible.

\begin{theorem}[Completion recovers the real Clifford algebras; Sylvester's law of inertia in GA]\label{c:synthesis.completion-recovers-the-real-clifford-algebras}
\changedmark{BurntOrange}{changed since the deposit}
Let $F$ be an ordered field. A quadratic space over $F$ is a finite-dimensional $F$-vector space $V$ with $Q(\lambda v) = \lambda^{2}Q(v)$ and polar form $B(u, v) = \tfrac12(Q(u + v) - Q(u) - Q(v))$ bilinear (as in \emph{Real quadratic space} [algebra core], with $F$ for $\R$); its Clifford algebra is defined by the universal property of \emph{The Clifford algebra of a quadratic space} [algebra core] over $F$, in which $uv + vu = 2B(u,v)$. $F$ is \emph{Euclidean} if every positive element is a square. $\Cl_{F}(p,q,r)$ is the Clifford algebra of $F^{n}$, $n = p + q + r$, with an orthogonal basis of $p$, $q$, $r$ vectors of squares $1$, $-1$, $0$.
\begin{enumerate}
\item[(i)] \emph{(Completeness gives square roots.)} If $F$ is complete (every nonempty subset bounded above has a least upper bound), then $F$ is Euclidean.
\item[(ii)] \emph{(Orthogonal bases by rejection.)} Every quadratic space over $F$ has a basis $u_{1}, \dots, u_{n}$ of mutually anticommuting vectors, $u_{i}u_{j} = -u_{j}u_{i}$ for $i \ne j$, obtained by repeated rejection $x \mapsto x - B(x, u)\,u^{-1}$ from vectors $u$ with $Q(u) \ne 0$. If $F$ is Euclidean, the basis can be chosen with $Q(e_{i}) \in \{1, -1, 0\}$, by $e_{i} = u_{i}/\sqrt{|Q(u_{i})|}$ where $Q(u_{i}) \ne 0$.
\item[(iii)] \emph{(Sylvester's law of inertia.)} For any ordered $F$, the numbers $p$, $q$, $r$ of vectors with $Q > 0$, $Q < 0$, $Q = 0$ in an anticommuting basis do not depend on the basis: $r = \dim \operatorname{rad} V$, and $p$ (respectively $q$) is the largest dimension of a subspace on which $Q$ is positive (negative) definite. If $F$ is Euclidean, two quadratic spaces are isometric if and only if they have the same $(p, q, r)$, and then $\Cl(V, Q) \cong \Cl_{F}(p,q,r)$.
\item[(iv)] \emph{(The completion recovers the store's algebra.)} Let $\Cl_{\Q}(p,q,r)$ be the $\Q$-span of the blades with the blade product of Proposition \ref{c:std-algebra.orthogonal-anticommute} (structure constants $0$, $\pm 1$). If $\R$ is a complete ordered field, as premise (F4) takes it, then $e_{S} \mapsto e_{S}$ extends to an isomorphism $\Cl_{\Q}(p,q,r) \otimes_{\Q} \R \cong \Cl(p,q,r)$, and for every real quadratic space $\Cl(V, Q) \cong \Cl_{\Q}(p,q,r) \otimes_{\Q} \R$ with $(p, q, r)$ its inertia; through $r$ added hyperbolic planes it embeds in the nondegenerate $\Cl_{\Q}(p{+}r,q{+}r) \otimes_{\Q} \R$, the radical generators going to null vectors $e'_{k} + f'_{k}$ (Theorem \ref{c:conformal.minimal-nondegenerate-completion}). So every real Clifford algebra of the store is, up to isomorphism, a rational Clifford algebra of a $\pm 1$ form, or a subalgebra of one, extended by scalars.
\item[(v)] \emph{(Over $\Q$ square classes remain.)} Over $\Q$, (ii) holds without the normalisation and the invariance of $(p, q, r)$ in (iii) holds, but the isometry criterion fails: $Q(\lambda u) = \lambda^{2}Q(u)$ fixes $Q(u)$ only up to a nonzero rational square. The line with $e^{2} = 2$ is not isometric to the line with $e^{2} = \pm 1$; its Clifford algebra, with $x\bar{x} = a^{2} - 2b^{2}$ for $x = a + be$, has no zero divisors, while $\Cl_{\Q}(1,0)$ has $(1 + e_{1})(1 - e_{1}) = 0$. The plane with $e_{1}^{2} = 2$, $e_{2}^{2} = 3$ has $x\bar{x} = a^{2} - 2b^{2} - 3c^{2} + 6d^{2}$ for $x = a + be_{1} + ce_{2} + de_{1}e_{2}$, which has no nontrivial rational zero, so its Clifford algebra is a division ring, while $\Cl_{\Q}(2,0)$ has zero divisors. Over any Euclidean field they become $\Cl(1,0)$ and $\Cl(2,0)$.
\item[(vi)] \emph{(What is used.)} Clauses (ii) to (iv) use of the coefficient field only that it is ordered and, for the normalisation, Euclidean. Completeness enters (iv) only through (i).
\end{enumerate}
\end{theorem}
\begin{proof}
The stored Propositions \ref{c:std-algebra.polarisation}, \ref{c:std-algebra.orthogonal-anticommute}, \ref{c:std-algebra.blade-basis} and \ref{c:std-algebra.functoriality}, and the existence of $\Cl(V, Q)$ as the quotient of the tensor algebra (\ref{c:std-algebra.existence-uniqueness}), are stated over $\R$; their proofs use only field operations and $\tfrac12$, and they are used here verbatim over $F$. Throughout, in a Clifford algebra over $F$, $(u + v)^{2} = Q(u + v)$ expands to $uv + vu = 2B(u, v)$, so two vectors anticommute exactly when $B(u, v) = 0$ (Proposition \ref{c:std-algebra.orthogonal-anticommute}, whose proof is the same over any field of characteristic $0$; an ordered field has characteristic $0$ since $m \cdot 1 > 0$ for $m \ge 1$). $B$ and $Q$ determine each other, $Q(v) = B(v, v)$ (Proposition \ref{c:std-algebra.polarisation}, again over any such field).

\emph{(i)} Let $a > 0$ and $S = \{x \in F : x \ge 0,\ x^{2} \le a\}$. $S$ contains $m = \min(1, a)$ (if $a \le 1$, $a^{2} \le a$; if $a > 1$, $1 \le a$), and is bounded above by $\max(1, a)$ (if $x > \max(1, a)$ then $x^{2} > x > a$). Let $s = \sup S \ge m > 0$. If $s^{2} < a$, put $h = \min\bigl(1, (a - s^{2})/(2s + 1)\bigr) > 0$; then $(s + h)^{2} = s^{2} + h(2s + h) \le s^{2} + h(2s + 1) \le a$, so $s + h \in S$, contradicting $s = \sup S$. If $s^{2} > a$, put $h = (s^{2} - a)/(2s)$, so $0 < h < s/2$; then $(s - h)^{2} = s^{2} - 2sh + h^{2} > s^{2} - 2sh = a$, and every $x \in S$ has $x \le s - h$ (otherwise $x^{2} > (s - h)^{2} > a$), so $s - h < s$ is an upper bound, contradicting minimality. Hence $s^{2} = a$. The check verifies both inequalities at rational instances, with the hand values $h = \tfrac13$, $(\tfrac43)^{2} = \tfrac{16}{9} \le 2$ and $h = \tfrac1{12}$, $(\tfrac{17}{12})^{2} = \tfrac{289}{144} > 2$ for $a = 2$, and that $2$ is not a rational square ($a^{2} \equiv 2b^{2} \bmod 4$ forces $a$, $b$ even).

\emph{(ii)} Induction on $n = \dim V$. If $B = 0$, every basis consists of mutually anticommuting null vectors. If $B \ne 0$ there is $u$ with $Q(u) \ne 0$: either some basis vector has $Q \ne 0$, or two have $B(x, y) \ne 0$ and then $Q(x + y) = 2B(x, y) \ne 0$. Such $u$ is invertible with $u^{-1} = u/Q(u)$ (since $uu = Q(u)$, a nonzero element of $F$; compare \emph{Which vectors and versors are invertible} [algebra core]), and the rejection $x_{\perp} = x - B(x, u)u^{-1} = \tfrac12(xu - ux)u^{-1}$ satisfies $B(x_{\perp}, u) = B(x, u) - B(x, u)B(u, u)/Q(u) = 0$, so $V = Fu \oplus u^{\perp}$, $x = B(x,u)u^{-1} + x_{\perp}$. By induction $u^{\perp}$ has an anticommuting basis, and with $u$ it is one for $V$, since every vector of $u^{\perp}$ anticommutes with $u$. If $F$ is Euclidean and $Q(u_{i}) \ne 0$, let $\lambda_{i} > 0$ with $\lambda_{i}^{2} = |Q(u_{i})|$; then $Q(u_{i}/\lambda_{i}) = Q(u_{i})/|Q(u_{i})| = \pm 1$, and rescaling keeps anticommutation. The check builds the rejected bases for six Gram matrices (hyperbolic, nondiagonal, degenerate) and verifies in the exact exterior model of the Clifford algebra that $u_{i}u_{j} + u_{j}u_{i} = 2B(u_{i}, u_{j})$ as operators, with $B(u_{i}, u_{j}) = 0$ for $i \ne j$; for $G = A^{T}DA$ with $A$ unitriangular and $D = \operatorname{diag}(4, -9, 0, \tfrac14)$ it recovers $D$ and the normalised squares $1, -1, 0, 1$.

\emph{(iii)} Let $u_{1}, \dots, u_{n}$ be an anticommuting basis, $P$, $N$, $Z$ the spans of the $u_{i}$ with $Q(u_{i}) > 0$, $< 0$, $= 0$, of dimensions $p$, $q$, $r$. For $v = \sum x_{i}u_{i}$, $B(v, u_{j}) = x_{j}Q(u_{j})$, so $v \in \operatorname{rad} V$ exactly when $x_{j} = 0$ for every $j$ with $Q(u_{j}) \ne 0$: $\operatorname{rad} V = Z$ and $r = \dim \operatorname{rad} V$, which does not depend on the basis. On $P$, $Q(\sum x_{i}u_{i}) = \sum x_{i}^{2}Q(u_{i}) > 0$ unless all $x_{i} = 0$, since in an ordered field a sum of positive multiples of squares, not all zero, is positive; so $P$ is a positive definite subspace of dimension $p$. On $N \oplus Z$, $Q \le 0$. If $W$ is positive definite with $\dim W > p$, then $\dim W + \dim(N \oplus Z) > n$; a basis of $W$ together with a basis of $N \oplus Z$ has more than $n$ vectors of $V$, so they are linearly dependent, and a dependence relation gives a nonzero $w \in W \cap (N \oplus Z)$, with $Q(w) > 0$ and $Q(w) \le 0$. So $p$ is the largest dimension of a positive definite subspace, and likewise $q$ for negative definite ones; both are independent of the basis. If $F$ is Euclidean, two spaces with the same $(p, q, r)$ have bases as in (ii) with the same squares; the linear map between them preserves $B$ on basis pairs, hence everywhere, and so preserves $Q$; conversely an isometry carries positive definite subspaces and radicals to the same, so the triples agree. By the universal property, an isometry $V \to F^{p,q,r}$ and its inverse extend to mutually inverse homomorphisms of the Clifford algebras (Proposition \ref{c:std-algebra.functoriality}, over $F$), so $\Cl(V, Q) \cong \Cl_{F}(p,q,r)$. The check orthogonalises six Gram matrices from four changes of basis and every ordering of the basis vectors, and finds one triple each, equal to the hand-stated signature, to $r = n - \operatorname{rank} G$, and to the congruence signature of GA.py.

\emph{(iv)} By Theorem \ref{c:std-algebra.blade-basis} the $2^{n}$ blades are an $\R$-basis of $\Cl(p,q,r)$, degenerate forms included, and by Proposition \ref{c:std-algebra.orthogonal-anticommute} $e_{S}e_{T} = \pm\bigl(\prod_{i \in S \cap T}Q(e_{i})\bigr)e_{S \triangle T}$ with the sign $(-1)^{\#\{(i,j) \in S \times T : i > j\}}$: the structure constants lie in $\{0, 1, -1\} \subset \Q$. They define $\Cl_{\Q}(p,q,r)$, and the $\R$-linear map $\Cl_{\Q}(p,q,r) \otimes_{\Q} \R \to \Cl(p,q,r)$, $e_{S} \otimes \lambda \mapsto \lambda e_{S}$, sends a basis to a basis and respects the products of basis elements, so it is an isomorphism. If $\R$ is complete, it is Euclidean by (i), so by (iii) every real quadratic space has some inertia $(p, q, r)$ and $\Cl(V, Q) \cong \Cl(p,q,r) \cong \Cl_{\Q}(p,q,r) \otimes_{\Q} \R$. Theorem \ref{c:conformal.minimal-nondegenerate-completion} embeds $\Cl(p,q,r)$ in $\Cl(p{+}r, q{+}r)$; explicitly, with $r$ added hyperbolic planes spanned by $e'_{k}$, $f'_{k}$ ($Q = 1, -1$), the vectors $e'_{k} + f'_{k}$ are null, anticommute with each other and with the old generators, so the universal property gives a homomorphism, rational on blades, whose images of the $2^{n}$ blades are linearly independent (each has a distinct leading blade). The check verifies that the structure constants of $\Cl(p,q,r)$ lie in $\{0, \pm 1\}$ for four signatures and that the null images realise $\Cl(p,q,r)$ injectively in $\Cl(p{+}r, q{+}r)$ (rank $2^{n}$) for five.

\emph{(v)} Over $\Q$, (ii) without the normalisation and the invariance part of (iii) hold as proved, since $\Q$ is ordered; the isometry criterion of (iii) used the normalisation and fails, as follows. If the line $\Q e$ with $Q(e) = 2$ were isometric to one with $Q(e') = \pm 1$, then $e \mapsto \lambda e'$ with $2 = \pm\lambda^{2}$, so $2$ would be a rational square; writing $\lambda = a/b$ in lowest terms, $a^{2} = 2b^{2}$ forces $a$ even, then $b$ even, a contradiction. In $\Cl(\langle 2 \rangle) = \Q \oplus \Q e$, $x\bar{x} = (a + be)(a - be) = a^{2} - 2b^{2}$, nonzero for $(a, b) \ne 0$ by the same argument, so $x^{-1} = \bar{x}/(x\bar{x})$; in $\Cl_{\Q}(1,0)$, $(1 + e_{1})(1 - e_{1}) = 0$. For the plane, $e_{1}$, $e_{2}$, $e_{1}e_{2}$ anticommute pairwise and $(e_{1}e_{2})^{2} = -e_{1}^{2}e_{2}^{2} = -6$, so for $x = a + be_{1} + ce_{2} + de_{1}e_{2}$ and its Clifford conjugate $\bar{x} = a - be_{1} - ce_{2} - de_{1}e_{2}$ the cross terms cancel and $x\bar{x} = a^{2} - 2b^{2} - 3c^{2} + 6d^{2}$. Suppose it vanishes for rationals not all zero; scale to integers with no common factor. Modulo $3$, $a^{2} + b^{2} \equiv 0$, and as squares are $0$ or $1$ modulo $3$, $3 \mid a$ and $3 \mid b$. Then modulo $9$, $-3(c^{2} - 2d^{2}) \equiv 0$, so $c^{2} + d^{2} \equiv 0 \bmod 3$ and $3 \mid c$, $3 \mid d$, a contradiction. So $x\bar{x} \ne 0$ for $x \ne 0$ and $x^{-1} = \bar{x}/(x\bar{x})$: a division ring. In $\Cl_{\Q}(2,0)$, $(1 + e_{1})(1 - e_{1}) = 0$. Over a Euclidean field, $e_{1}/\sqrt{2}$ and $e_{2}/\sqrt{3}$ are anticommuting vectors of square $1$, so the line and the plane become $\Cl(1,0)$ and $\Cl(2,0)$. The check computes both norms $x\bar{x}$ with GA.py on a grid, confirms that they vanish only at $0$ there and give inverses, runs the two descent lemmas over all residues modulo $4$ and modulo $9$, and exhibits the zero divisors of $\Cl_{\Q}(1,0)$ and $\Cl_{\Q}(2,0)$.

\emph{(vi)} Inspection: (ii) uses only that $Q(u) \ne 0$ can be divided by, and square roots of positive elements for the normalisation; (iii) uses the order (sums of positive multiples of squares are positive) and, for the isometry statement, the normalisation; (iv) uses completeness only to invoke (i).

\end{proof}

\phantomsection\label{c:note.geometry-first.setting-completion-number-systems-inside-2}
\noindent\emph{How the proof goes.} Square roots come from the least upper bound of the numbers whose square is at most $a$; orthogonal bases come by induction, splitting off one vector of nonzero square at a time; the counts are then compared directly. These are the classical arguments, written out so that they use (F4) and nothing more. The proof is in Appendix~\ref{prf:thm:bitinertia}.

So over $\R$ a quadratic form is known by three counts, and when it is nondegenerate its algebra is one of the algebras $A_{\R} = \Cl(p,q)$ the bit produced. Part (iii) used only the order and part (ii) only square roots of positive numbers. Over $\Q$, where those roots are missing, a direction carries more than a sign: its square is fixed only up to multiplication by a nonzero rational square, so that $x^{2}$ and $2x^{2}$ are different forms with the same counts.

The second stage needs two standard tools of polynomial algebra, proved here for completeness.

\phantomsection\label{c:note.geometry-first.setting-completion-number-systems-inside-3}
\noindent\emph{How the proof goes.} Both are classical: (a) by adjoining a root of an irreducible factor and inducting on the degree, (b) by removing the leading term again and again. The proof is in Appendix~\ref{prf:lem:bitpoly}.

Write $\C$ for the algebra $A_{\R}$ with one direction $e$ of square $-1$, that is $\Cl(0,1)$: its elements are $z = a + be$ with $a, b \in \R$. It is commutative, and with $\bar{z} = a - be$,
\[
z\,\bar{z} \;=\; a^{2} + b^{2} ,
\]
which is positive unless $z = 0$. So every nonzero $z$ has the inverse $\bar{z}/(z\bar{z})$: $\C$ is a field, the map $z \mapsto \bar{z}$ respects sums and products, and the elements it fixes are those of $\R$.

\phantomsection\label{c:note.geometry-first.setting-completion-number-systems-inside-4}
\noindent\emph{How the proof goes.} (i) uses the least upper bound, (ii) writes the square roots down, and (iii) is Laplace's argument, an induction on the power of $2$ in the degree that uses Lemma~\ref{lem:bitpoly}; (iv) follows from (iii). Example~\ref{c:note.geometry-first.roots-cl-0-1-hamiltons-relations} shows a pair of roots. The proof is in Appendix~\ref{prf:thm:bitroots}.

This is the fundamental theorem of algebra, by Laplace's argument \cite{Ebbinghaus1991}. (The same holds for a nonconstant polynomial $f$ with coefficients in $\C$: the product of $f$ and the polynomial with the conjugate coefficients has real coefficients, so it has a root $z \in \C$, and then $f(z) = 0$ or $f(\bar{z}) = 0$.) The theorem places the complex numbers: they are not adjoined to the real numbers from outside but are the algebra of one direction of square $-1$, and that algebra already holds the roots of every real polynomial. Completeness was used twice only, for square roots and for (i).

\begin{theorem}[Frobenius in GA; every finite-dimensional associative real algebra in a corner of Cl(k,k)]\label{c:synthesis.frobenius-and-universality-in-clifford-algebras}
\changedmark{BurntOrange}{changed since the deposit}
Let $\R$ be a complete ordered field, as premise (F4) takes the coefficients to be. $\Cl(0,1) = \mathrm{span}\{1, e_{1}\}$, $e_{1}^{2} = -1$, is the complex field $\C$ (\ref{c:std-algebra.complex-numbers}), and $\Cl(0,2)$ is the quaternions $\HH$ (\ref{c:std-algebra.quaternions}).
\begin{enumerate}
\item[(i)] \emph{(Real polynomials have their roots in $\Cl(0,1)$.)} Every nonconstant polynomial with coefficients in $\R$ has a root in $\Cl(0,1)$. Hence every irreducible polynomial over $\R$ has degree $1$ or $2$. (Completeness is used only for square roots of positive elements and a root of each odd-degree polynomial, so (i) and (ii) hold over any real closed field.)
\item[(ii)] \emph{(Frobenius, in geometric algebra.)} Let $D \ne 0$ be a finite-dimensional associative unital $\R$-algebra without zero divisors. Then $V = \{x \in D : x^{2} \in \R,\ x^{2} \le 0\}$ is a subspace, $D = \R 1 \oplus V$, and $Q(v) = v^{2}$ is a negative definite quadratic form on $V$ with polar form $B(u, v) = \tfrac12(uv + vu)$. The inclusion $V \subseteq D$ extends to a surjective homomorphism $\Cl(0,k) \to D$, $k = \dim V$. If $k \ge 2$, the bivector $e_{1}e_{2}$ of an orthonormal basis lies in $V$, is orthogonal to $e_{1}$ and $e_{2}$, and commutes with every $e_{j}$, $j \ge 3$, which forces $e_{j} = \pm e_{1}e_{2}$; so $k \in \{0, 1, 3\}$ and $D$ is isomorphic to $\R$, $\Cl(0,1)$ or $\Cl(0,2)$ (for $k = 3$, $e_{3} = \pm e_{1}e_{2}$: the quaternionic quotient of $\Cl(0,3)$). Conversely these three have no zero divisors (\ref{c:transport.exceptional-division-algebras}).
\item[(iii)] \emph{(Every finite-dimensional real algebra sits in some $\Cl(k,k)$.)} Let $v_{i}$, $w_{i}$ ($1 \le i \le k$) be a Witt basis of $\R^{k,k}$, $f = v_{1}w_{1}v_{2}w_{2}\cdots v_{k}w_{k}$, and for $I, J \subseteq \{1, \dots, k\}$ let $E_{I,J} = w_{I}\,f\,v_{J}^{\leftarrow}$, with $w_{I}$ the increasing and $v_{J}^{\leftarrow}$ the decreasing product. The $E_{I,J}$ are a complete system of $2^{k} \times 2^{k}$ matrix units of $\Cl(k,k)$. If $A$ is an associative unital $\R$-algebra of finite dimension $d \le 2^{k}$, $L(a)$ the matrix of left multiplication by $a$ in a basis of $A$, and $I_{1}, \dots, I_{d}$ distinct index sets, then $\varphi(a) = \sum_{s,t}L(a)_{st}E_{I_{s},I_{t}}$ is an injective algebra homomorphism $A \to \Cl(k,k)$ whose image of $1$ is the idempotent $\epsilon = \sum_{s}E_{I_{s},I_{s}}$: $A$ is a subalgebra of the corner $\epsilon\,\Cl(k,k)\,\epsilon \cong \Mat_{d}(\R)$.
\item[(iv)] \emph{(The number systems inside.)} $\R = \Cl(0,0)$; $\C = \Cl(0,1)$, also the even part of $\Cl(2,0)$ (\ref{c:std-algebra.complex-numbers}); the split-complex numbers $\Cl(1,0) = \R\epsilon_{+} \oplus \R\epsilon_{-}$, $\epsilon_{\pm} = \tfrac12(1 \pm e_{1})$, $\cong \R \oplus \R$; the dual numbers $\mathrm{span}\{1, n\} \cong \Cl(0,0,1)$ for a nonzero null vector $n$ of $\Cl(1,1)$; $\HH = \Cl(0,2)$, also the even part of $\Cl(3,0)$; and $\Mat_{2^{k}}(\R) \cong \Cl(k,k)$. By (ii) no other finite-dimensional associative real division algebra exists, and by (iii) every finite-dimensional associative real algebra is a subalgebra of a corner $\epsilon\,\Cl(k,k)\,\epsilon$ of some $\Cl(k,k)$ (in general not a unital subalgebra of $\Cl(k,k)$: $\Mat_{3}(\R)$ has no unital embedding in any $\Mat_{2^{k}}(\R)$).
\end{enumerate}
\end{theorem}
\begin{proof}
Write $\C$ for $\Cl(0,1) = \mathrm{span}\{1, e_{1}\}$, a commutative field (\ref{c:std-algebra.complex-numbers}), with conjugation $\overline{a + be_{1}} = a - be_{1}$, a ring automorphism fixing $\R$, and $z\bar{z} = a^{2} + b^{2}$.

\emph{(i), step 1: odd degree.} Let $p = X^{n} + \sum_{i < n}c_{i}X^{i}$, $n$ odd, and $M = 1 + \sum_{i < n}|c_{i}|$. For $|x| \ge M \ge 1$, $|\sum_{i<n}c_{i}x^{i}| \le (M - 1)|x|^{n-1} < |x|^{n}$, so $p(M) > 0 > p(-M)$. On $[-M, M]$, $|p(x) - p(y)| \le \Lambda|x - y|$ with $\Lambda = \sum_{l}l\,|c_{l}|M^{l-1}$ ($c_{n} = 1$), since $x^{i} - y^{i} = (x - y)\sum_{j}x^{j}y^{i-1-j}$. Let $s = \sup\{x \in [-M, M] : p(x) < 0\}$, which exists by completeness. If $p(s) < 0$, then $s < M$ and $p < 0$ on $[s, s + |p(s)|/\Lambda)$, against $s$ being an upper bound; if $p(s) > 0$, then $p > 0$ on $(s - p(s)/\Lambda, s]$, so $s - p(s)/\Lambda$ is a smaller upper bound. Hence $p(s) = 0$. The same argument applied to $X^{2} - a$, $a > 0$, on $[0, M]$ with $M = 1 + a$ (where $p(0) = -a < 0 < p(M)$) gives $\sqrt{a} \in \R$.

\emph{Step 2: square roots in $\C$.} For $z = a + be_{1}$ let $r = \sqrt{a^{2} + b^{2}} \ge |a|$, $x = \sqrt{(r + a)/2}$, $y = \sigma\sqrt{(r - a)/2}$ with $\sigma = 1$ if $b \ge 0$ and $-1$ otherwise. Then $(x + ye_{1})^{2} = x^{2} - y^{2} + 2xye_{1} = a + \sigma\sqrt{r^{2} - a^{2}}\,e_{1} = a + be_{1}$. So every monic quadratic $X^{2} - sX + P$ over $\C$ factors as $(X - \tfrac12(s + w))(X - \tfrac12(s - w))$ with $w^{2} = s^{2} - 4P$.

\emph{Step 3: Laplace's induction.} Write $n = \deg p = 2^{m}u$ with $u$ odd; we show by induction on $m$ that $p$ has a root in $\C$. For $m = 0$ this is step 1. For $m \ge 1$, let $L$ be a splitting field of $p$ over the field $\C = \Cl(0,1)$ (obtained by adjoining roots of irreducible factors one at a time, $K[X]/(g)$ being a field for $g$ irreducible over a field $K$), so that $\R \subseteq \C \subseteq L$. Let $r_{1}, \dots, r_{n} \in L$ be the roots of $p$ with multiplicity. For $t \in \R$ put $g_{t}(X) = \prod_{i<j}(X - (r_{i} + r_{j} + t\,r_{i}r_{j}))$, of degree $n(n-1)/2 = 2^{m-1}u(n - 1)$ with $u(n - 1)$ odd. Its coefficients are polynomials over $\R$ in the $r_{i}$, invariant under permutations of the $r_{i}$, hence (fundamental theorem on symmetric polynomials: by induction on the leading monomial in lexicographic order, such a polynomial is a polynomial in the elementary symmetric ones) polynomials over $\R$ in the coefficients of $p$: $g_{t}$ has real coefficients. By induction it has a root in $\C$, so for each $t$ some pair $i < j$ has $r_{i} + r_{j} + t\,r_{i}r_{j} \in \C$. There are infinitely many $t$ and finitely many pairs, so one pair serves two values $t \ne t'$; subtracting, $(t - t')r_{i}r_{j} \in \C$, so $P = r_{i}r_{j} \in \C$ and $s = r_{i} + r_{j} \in \C$. Then $r_{i}$ is a root of $X^{2} - sX + P$, whose roots in the field $L$ are those of step 2, which lie in $\C$. Finally, if $q$ is irreducible over $\R$ of degree at least $2$, it has a root $z \in \C \setminus \R$, also $\bar{z}$ is a root, and $(X - z)(X - \bar{z}) = X^{2} - (z + \bar{z})X + z\bar{z}$ has real coefficients and divides $q$ over $\C$, hence over $\R$ (division with remainder over $\R$ gives the same quotient and remainder as over $\C$); so $\deg q = 2$. The check verifies the square-root formula on squares of rational elements of $\Cl(0,1)$, the sign change of three odd-degree polynomials at $\pm M$, and for $p = (X^{2} + 1)(X^{2} + 4)$, with roots $\pm e_{1}$, $\pm 2e_{1}$ in $\Cl(0,1)$, that $g_{t}$ has real coefficients for five values of $t$ and equals $X^{6} + 10X^{4} + 9X^{2}$ at $t = 0$.

\emph{(ii)} A nonzero $x \in D$ is invertible: $y \mapsto xy$ is injective and linear on the finite-dimensional $D$, hence onto, so $xy = 1$ for some $y$, and $x(yx - 1) = 0$ gives $yx = 1$. For $x \in D$ the powers $1, x, \dots, x^{\dim D}$ are dependent, so $x$ has a monic minimal polynomial $m_{x}$ over $\R$; it is irreducible, since $m_{x} = gh$ with smaller degrees would give $g(x)h(x) = 0$ with both factors nonzero. By (i) its degree is $1$ or $2$. If $x \notin \R 1$, $m_{x} = X^{2} - 2bX + c$ with $b^{2} - c < 0$ (no real root), and $(x - b)^{2} = b^{2} - c < 0$, so $x - b \in V$: $D = \R 1 + V$. A real number with non-positive square is $0$, so $\R 1 \cap V = 0$. For $v \in V \setminus 0$, $v^{2} < 0$ (as $v^{2} = 0$ would make $v$ a zero divisor) and $v \notin \R 1$. $V$ is closed under real multiples, $(\lambda v)^{2} = \lambda^{2}v^{2}$. For $u, v \in V$: if $u = 0$ there is nothing to show; if $1, u, v$ are dependent and $u \ne 0$, then $v = \alpha + \beta u$ and $v^{2} = \alpha^{2} + \beta^{2}u^{2} + 2\alpha\beta u$ is real, so $\alpha\beta = 0$; $\beta = 0$ gives $v = \alpha \in \R 1 \cap V = 0$, and otherwise $u + v = (1 + \beta)u \in V$. If $1, u, v$ are independent, $u \pm v \notin \R 1$, so $(u + v)^{2} = \alpha(u + v) + \beta$ and $(u - v)^{2} = \gamma(u - v) + \delta$; adding, $2u^{2} + 2v^{2} = (\alpha + \gamma)u + (\alpha - \gamma)v + \beta + \delta$, and independence gives $\alpha = \gamma = 0$, so $(u + v)^{2} = \beta \in \R$. If $\beta > 0$, $(u + v - \sqrt{\beta})(u + v + \sqrt{\beta}) = 0$ would put $u + v \in \R 1$; $\beta = 0$ would give $u + v = 0$; so $\beta < 0$ and $u + v \in V$. Thus $V$ is a subspace and $D = \R 1 \oplus V$. For $u, v \in V$, $uv + vu = (u + v)^{2} - u^{2} - v^{2} \in \R$, so $B(u, v) = \tfrac12(uv + vu)$ is a symmetric bilinear form with $B(v, v) = v^{2} = Q(v) < 0$ for $v \ne 0$. Rejection $x \mapsto x - B(x, u)u^{-1}$ and division by $\sqrt{-Q(u)}$ (step 1) give an orthonormal basis $e_{1}, \dots, e_{k}$ of $V$: $e_{i}^{2} = -1$, $e_{i}e_{j} = -e_{j}e_{i}$ for $i \ne j$ (Proposition \ref{c:std-algebra.orthogonal-anticommute}). By the universal property the inclusion of $V$ extends to a homomorphism $\phi: \Cl(0,k) \to D$, onto since its image contains $1$ and $V$; $\dim \Cl(0,k) = 2^{k}$ (Theorem \ref{c:std-algebra.blade-basis}).

Let $k \ge 2$ and $u = e_{1}e_{2} \in D$. Then $u^{2} = -e_{1}^{2}e_{2}^{2} = -1$, so $u \in V$, and $ue_{1} + e_{1}u = e_{1}e_{2}e_{1} + e_{1}e_{1}e_{2} = e_{2} - e_{2} = 0$, likewise $ue_{2} + e_{2}u = 0$: $B(u, e_{1}) = B(u, e_{2}) = 0$. Writing $u = \sum_{i}c_{i}e_{i}$, $B(u, e_{j}) = -c_{j}$. If $k = 2$ this gives $u = 0$, against $u^{2} = -1$. If $k \ge 3$ and $j \ge 3$, $e_{j}$ anticommutes with $e_{1}$ and $e_{2}$, hence commutes with $u$, and $2ue_{j} = ue_{j} + e_{j}u = 2B(u, e_{j}) = -2c_{j}$ is real, so $e_{j} = u^{-1}(ue_{j}) = c_{j}u$ (as $u^{-1} = -u$). The vectors $e_{3}, \dots, e_{k}$ are independent and all lie in $\R u$, so $k \le 3$, and for $k = 3$, $-1 = e_{3}^{2} = -c_{3}^{2}$ gives $e_{3} = \pm e_{1}e_{2}$. Hence: $k = 0$ and $D = \R$; or $k = 1$ and $\phi: \Cl(0,1) \to D$ is onto between spaces of dimension $2$, an isomorphism; or $k = 3$, $D = \mathrm{span}\{1, e_{1}, e_{2}, e_{1}e_{2}\}$ is the image of the subalgebra generated by $e_{1}, e_{2}$, a copy of $\Cl(0,2)$ of dimension $4$, so $D \cong \Cl(0,2) \cong \HH$ (\ref{c:std-algebra.quaternions}). In $\Cl(0,3)$ itself the element $ue_{3} = e_{1}e_{2}e_{3}$ is central with square $+1$, and $(1 - ue_{3})(1 + ue_{3}) = 0$: the map to $D$ sends it to $\pm 1$, the quaternionic quotient. That $\R$, $\Cl(0,1)$, $\Cl(0,2)$ have no zero divisors is Theorem \ref{c:transport.exceptional-division-algebras}. The check verifies in $\Cl(0,1)$ and $\Cl(0,2)$ that an element squares to a real exactly when its real or its pure part vanishes and that $x^{2} - 2ax + (a^{2} - v^{2}) = 0$; in $\Cl(0,k)$, $k = 2, 3, 4$, the relations of $u = e_{1}e_{2}$ with the generators and $(1 - ue_{j})(1 + ue_{j}) = 0$; that $e_{1}, e_{2}, e_{1}e_{2}$ satisfy the relations of three generators of $\Cl(0,3)$; that $e_{1}e_{2} \notin \mathrm{span}\{1, e_{1}, e_{2}\}$; and the dimensions $1, 2, 4$.

\emph{(iii)} By Proposition \ref{c:std-algebra.witt-basis-car} the $v_{i}$, $w_{i}$ satisfy $v_{i}^{2} = w_{i}^{2} = 0$, $v_{i}w_{i} + w_{i}v_{i} = 1$, and anticommute across different indices, and the $v_{i}w_{i}$ are commuting idempotents, so $f$ is idempotent. Each $v_{i}w_{i}$ commutes with $v_{j}$, $w_{j}$ for $j \ne i$ (two sign changes), so $v_{i}f = 0$ and $fw_{i} = 0$ (move the factor $v_{i}w_{i}$ next to $v_{i}$, respectively $w_{i}$, and use $v_{i}^{2} = w_{i}^{2} = 0$). Consider $fv_{J}^{\leftarrow}w_{K}f$. If $i \in K \setminus J$, move $w_{i}$ to the left through the $v_{j}$, $j \in J$, and the $w$'s before it, all anticommuting with it, to meet $f$: the product is $0$. If $i \in J \setminus K$, move $v_{i}$ to the right likewise: $0$. If $J = K$, the innermost pair $v_{j}w_{j}$ is even and commutes with the remaining letters, so $v_{J}^{\leftarrow}w_{J} = \prod_{j \in J}v_{j}w_{j}$, and $\prod_{j \in J}(v_{j}w_{j})f = f$. So $fv_{J}^{\leftarrow}w_{K}f = \delta_{JK}f$ and $E_{I,J}E_{K,L} = \delta_{JK}E_{I,L}$. By the same moves $E_{I,I} = \prod_{i \in I}w_{i}v_{i}\prod_{i \notin I}v_{i}w_{i}$, and $\sum_{I}E_{I,I} = \prod_{i}(w_{i}v_{i} + v_{i}w_{i}) = 1$. So the $E_{I,J}$ are $2^{k} \times 2^{k}$ matrix units; by Proposition \ref{c:std-algebra.matrix-units} $\Cl(k,k) \cong \Mat_{2^{k}}(f\Cl(k,k)f)$, and comparing dimensions ($4^{k}$, Theorem \ref{c:std-algebra.fock-representation}) the corner $f\Cl(k,k)f$ is $\R f$. Now $a \mapsto L(a)$ is an injective unital homomorphism $A \to \Mat_{d}(\R)$ ($L(a)$ applied to the coordinates of $1$ gives those of $a$), and $M \mapsto \sum_{s,t}M_{st}E_{I_{s},I_{t}}$ is an injective homomorphism $\Mat_{d}(\R) \to \Cl(k,k)$ (the $E$'s multiply as matrix units and are linearly independent) sending the identity to $\epsilon$; the composite is $\varphi$, and $\varphi(A) = \epsilon\varphi(A)\epsilon$. The check verifies the matrix-unit relations for $k = 1, 2$, $\sum_{I}E_{I,I} = 1$ for $k = 1, 2, 3$, and the embedding (homomorphism on basis pairs, injectivity, the idempotent unit) for the dual numbers, the split-complex numbers, $\R[x]/(x^{3})$, the upper triangular $2 \times 2$ matrices and $\Cl(0,2)$.

\emph{(iv)} $\epsilon_{\pm}^{2} = \tfrac14(1 \pm 2e_{1} + e_{1}^{2}) = \epsilon_{\pm}$, $\epsilon_{+}\epsilon_{-} = \tfrac14(1 - e_{1}^{2}) = 0$, $\epsilon_{+} + \epsilon_{-} = 1$, so $\Cl(1,0) = \R\epsilon_{+} \oplus \R\epsilon_{-} \cong \R \oplus \R$. In $\Cl(1,1)$, $n = e_{1} + e_{2}$ has $n^{2} = 1 - 1 = 0$, so $\mathrm{span}\{1, n\}$ is closed, $(a + bn)(c + dn) = ac + (ad + bc)n$, and $\rho \mapsto n$ is an isomorphism from $\Cl(0,0,1)$. The complex numbers and quaternions are Propositions \ref{c:std-algebra.complex-numbers} and \ref{c:std-algebra.quaternions}, and $\Cl(k,k) \cong \Mat_{2^{k}}(\R)$ is (iii) with Theorem \ref{c:std-algebra.fock-representation}. The final sentence restates (ii) and (iii). The check verifies the idempotents of $\Cl(1,0)$, the dual-number product on a null vector of $\Cl(1,1)$, and $x\bar{x} = $ sum of squares in $\Cl(0,1)$ and $\Cl(0,2)$.

\end{proof}

\phantomsection\label{c:note.geometry-first.setting-completion-number-systems-inside-5}
\noindent\emph{How the proof goes.} Every element satisfies a real polynomial of degree at most $2$; the elements whose square is real and not positive form a subspace $V$, on which $-x^{2}$ is a positive form; $V$ has dimension $0$, $1$ or $3$, and in the last case the third direction is the product of the first two (Palais's argument). The proof is in Appendix~\ref{prf:thm:bitfrobenius}.

The theorem is due to Frobenius \cite{Frobenius1878}; the proof through the subspace of elements with non-positive square follows Palais \cite{Palais1968}. Read in the algebra, it says that the only directions that can be added without ever losing division are at most two, both of square $-1$: a third must be the product of the first two.

With these three theorems the order of the textbooks is reversed. The real numbers are the scalars of the completed algebra; the complex numbers and the quaternions are two of its smallest members; and by the classification of the real Clifford algebras (imported from \cite[Theorem 9.3]{ButlerT0}, where its induction is cited and not repeated, and used here in this sentence only, in no proof), every other $\Cl(p,q)$ is a matrix algebra over one of the three, or a sum of two. The number systems in common use sit inside the algebra the bit generates.

\begin{example}[Roots in $\Cl(0,1)$, and Hamilton's relations]\label{c:note.geometry-first.roots-cl-0-1-hamiltons-relations}
\changedmark{BurntOrange}{changed since the deposit}
The polynomial $x^{2} - 2x + 5$ has no real root, and in $\C = \Cl(0,1)$, with $e^{2} = -1$, it has the roots $1 \pm 2e$:
\[
(1 + 2e)^{2} - 2(1 + 2e) + 5 \;=\; (1 + 4e - 4) - 2 - 4e + 5 \;=\; 0 ,
\]
and likewise for $1 - 2e$. In $\HH = \Cl(0,2)$ put $\mathbf{i} = e_{1}$, $\mathbf{j} = e_{2}$ and $\mathbf{k} = e_{1}e_{2}$. Then
\[
\mathbf{i}^{2} = \mathbf{j}^{2} = \mathbf{k}^{2} = -1, \qquad \mathbf{i}\mathbf{j} = \mathbf{k} = -\mathbf{j}\mathbf{i} ,
\]
since $\mathbf{k}^{2} = e_{1}e_{2}e_{1}e_{2} = -e_{1}e_{1}e_{2}e_{2} = -1$: these are Hamilton's relations. A third direction of square $-1$ breaks division, as Theorem~\ref{c:synthesis.frobenius-and-universality-in-clifford-algebras} says it must. In $\Cl(0,3)$ the element $u = e_{1}e_{2}e_{3}$ has $u^{2} = 1$, and so
\[
(1 - u)(1 + u) \;=\; 1 - u^{2} \;=\; 0
\]
with both factors nonzero.
\end{example}

\phantomsection\label{c:note.geometry-first.setting-completion-number-systems-inside-6}
Two remarks on what (F4) is needed for. Theorems \ref{thm:bitroots} and \ref{c:synthesis.frobenius-and-universality-in-clifford-algebras} use completeness only through square roots of positive numbers and roots of polynomials of odd degree, so they hold over every ordered field with those two properties, for example the real algebraic numbers. Completeness itself is needed for limits, such as the half-angle rotor of \cite[Example 6.1]{ButlerT0}, which needs completeness to be defined for every angle.

One family is not reached by any finite step. The commutation relation of a boson, $ab - ba = 1$, has no solution in a nonzero finite-dimensional real algebra: acting on the algebra by left multiplication, $ab - ba$ has trace $0$ and $1$ has trace the dimension. Whether such relations arise as a limit on the infinite tower of T3 \cite{ButlerT3} lies outside this note.

\section{The switch: what a bit does}

\phantomsection\label{c:note.geometry-first.setting-switch-what-bit-does}
The ladder above treats a distinction as something that can be made and made again, and the direction $e_{i}$ is that operation. A bit can also be acted on in a less symmetric way: it can be \emph{set} and it can be \emph{cleared}. This section shows that those two operations are exactly a hyperbolic pair, that they need no fractions, and that the integer algebra of the ladder does not contain them.

\begin{theorem}[Set and clear of a bit generate the hyperbolic plane]\label{c:synthesis.set-and-clear-of-a-bit-generate-the-hyperbolic-plane}
\changedmark{BurntOrange}{changed since the deposit}
Let $R$ be a commutative ring. A \emph{bit over $R$} is the free module $M = R\,s_{\mathrm{off}} \oplus R\,s_{\mathrm{on}}$ on two states; \emph{set} is the $R$-linear map $v$ with $v\,s_{\mathrm{off}} = s_{\mathrm{on}}$, $v\,s_{\mathrm{on}} = 0$, and \emph{clear} the $R$-linear map $w$ with $w\,s_{\mathrm{on}} = s_{\mathrm{off}}$, $w\,s_{\mathrm{off}} = 0$. A \emph{switch} in an associative unital ring is a pair $(a, b)$ with $a^{2} = b^{2} = 0$ and $ab + ba = 1$.
\begin{enumerate}
\item[(i)] $(v, w)$ is a switch. The products $vw$ and $wv$ are idempotents with $vw + wv = 1$ and $vw \cdot wv = wv \cdot vw = 0$; $vw$ fixes $s_{\mathrm{on}}$ and annihilates $s_{\mathrm{off}}$ (the test ``is on''), and $wv$ the reverse.
\item[(ii)] The $R$-algebra generated by $v$ and $w$ is $\mathrm{End}_{R}(M)$, free with basis $1, v, w, vw$.
\item[(iii)] For $x = \alpha v + \beta w$ with $\alpha, \beta \in R$, $x^{2} = \alpha\beta$. With $L = Rv \oplus Rw$ and $Q(x) = \alpha\beta$ (the hyperbolic quadratic module): for every associative unital $R$-algebra $B$ and $R$-linear $\phi : L \to B$ with $\phi(x)^{2} = Q(x)\,1$ there is exactly one unital $R$-algebra homomorphism $\Phi : \mathrm{End}_{R}(M) \to B$ extending $\phi$. So $\mathrm{End}_{R}(M)$ is the Clifford algebra of the hyperbolic plane over every $R$, with no hypothesis on $2$.
\item[(iv)] $e_{+} = v + w$ and $e_{-} = v - w$ satisfy $e_{+}^{2} = 1$, $e_{-}^{2} = -1$, $e_{+}e_{-} = -e_{-}e_{+} = wv - vw$; $e_{+}$ exchanges the two states and $e_{+}e_{-}$ is $+1$ on $s_{\mathrm{off}}$ and $-1$ on $s_{\mathrm{on}}$.
\end{enumerate}
\end{theorem}
\begin{proof}
\emph{(i)} $v^{2}$ sends $s_{\mathrm{off}}$ to $v\,s_{\mathrm{on}} = 0$ and $s_{\mathrm{on}}$ to $0$; likewise $w^{2} = 0$. $vw$ sends $s_{\mathrm{on}}$ to $v\,s_{\mathrm{off}} = s_{\mathrm{on}}$ and $s_{\mathrm{off}}$ to $0$; $wv$ sends $s_{\mathrm{off}}$ to $s_{\mathrm{off}}$ and $s_{\mathrm{on}}$ to $0$. So $vw$ and $wv$ are the two coordinate projections: idempotent, with sum $1$ and products $0$.

\emph{(ii)} In the basis $(s_{\mathrm{off}}, s_{\mathrm{on}})$ the maps $wv$, $w$, $v$, $vw$ are the four matrices with a single entry $1$, a basis of the $2 \times 2$ matrices over $R$, and $wv = 1 - vw$; so $1, v, w, vw$ is a basis as well (the change of basis has determinant $\pm 1$).

\emph{(iii)} $x^{2} = \alpha^{2}v^{2} + \beta^{2}w^{2} + \alpha\beta\,(vw + wv) = \alpha\beta$. Let $\phi : L \to B$ be $R$-linear with $\phi(x)^{2} = Q(x)\,1$, and put $a = \phi(v)$, $b = \phi(w)$. Then $a^{2} = Q(v) = 0$, $b^{2} = Q(w) = 0$ and $ab + ba = (a + b)^{2} - a^{2} - b^{2} = Q(v + w) = 1$: $(a, b)$ is a switch in $B$. Define the $R$-linear map $\Phi$ on the basis of (ii) by $1 \mapsto 1$, $v \mapsto a$, $w \mapsto b$, $vw \mapsto ab$. The sixteen products of basis elements are reduced to the basis using only $v^{2} = w^{2} = 0$ and $wv = 1 - vw$: $v \cdot w = vw$, $w \cdot v = 1 - vw$, $v \cdot vw = 0$, $vw \cdot v = v(1 - vw) = v$, $w \cdot vw = (1 - vw)w = w$, $vw \cdot w = 0$, $vw \cdot vw = v(1 - vw)w = vw$, and the rest are $0$ or involve $1$. The same reductions hold for $a$, $b$, which satisfy the same relations, so $\Phi$ is multiplicative on the basis and hence a homomorphism. It is unique because $v$ and $w$ generate $\mathrm{End}_{R}(M)$ by (ii).

\emph{(iv)} $(v \pm w)^{2} = v^{2} + w^{2} \pm (vw + wv) = \pm 1$, and $(v + w)(v - w) = v^{2} - vw + wv - w^{2} = wv - vw = -(v - w)(v + w)$. By (i), $v + w$ sends $s_{\mathrm{off}}$ to $s_{\mathrm{on}}$ and back, and $wv - vw$ is $+1$ on $s_{\mathrm{off}}$ and $-1$ on $s_{\mathrm{on}}$.

\end{proof}

\phantomsection\label{c:note.geometry-first.setting-switch-what-bit-does-2}
\noindent\emph{How the proof goes.} Everything is read off from what $v$ and $w$ do to the two states. Example~\ref{c:note.geometry-first.switch-two-two-tables} shows each identity as a two-by-two table. The proof is in Appendix~\ref{prf:thm:switch}.

Nothing in the statement or the proof divides by $2$, and no sign was put in: the relations of a single fermionic mode, $v^{2} = w^{2} = 0$ and $vw + wv = 1$ \cite{JordanWigner1928}, are the relations of setting and clearing a bit. ``It cannot be set twice'' is $v^{2} = 0$, and ``it is either on or off'' is $vw + wv = 1$. That the Clifford algebra of the hyperbolic plane is the $2 \times 2$ matrix ring over every commutative ring is classical \cite{Bass1974, Knus1991}. The primer's Example 4.1 \cite{ButlerT0} met the same pair from the other side, as $\tfrac12(e_{1} \pm e_{2})$, where the $\tfrac12$ was needed; part (iv) is that example read backwards, and it produces the two directions of squares $+1$ and $-1$ from the switch without any fraction.

\begin{example}[The switch as two-by-two tables]\label{c:note.geometry-first.switch-two-two-tables}
In the basis $s_{\mathrm{off}}, s_{\mathrm{on}}$, each column giving the image of a state,
\[
v = \begin{pmatrix} 0 & 0 \\ 1 & 0 \end{pmatrix}, \qquad
w = \begin{pmatrix} 0 & 1 \\ 0 & 0 \end{pmatrix}, \qquad
vw = \begin{pmatrix} 0 & 0 \\ 0 & 1 \end{pmatrix}, \qquad
wv = \begin{pmatrix} 1 & 0 \\ 0 & 0 \end{pmatrix} .
\]
So $v^{2} = w^{2} = 0$ (setting twice gives nothing), $vw$ is the test for on and $wv$ the test for off, and $vw + wv = 1$. The sum and the difference are
\[
e_{+} = v + w = \begin{pmatrix} 0 & 1 \\ 1 & 0 \end{pmatrix}, \qquad
e_{-} = v - w = \begin{pmatrix} 0 & -1 \\ 1 & 0 \end{pmatrix},
\]
with $e_{+}^{2} = 1$ (exchanging the states twice changes nothing), $e_{-}^{2} = -1$, and $e_{+}e_{-} = \begin{pmatrix} 1 & 0 \\ 0 & -1 \end{pmatrix}$. No entry is a fraction.
\end{example}

\phantomsection\label{c:note.geometry-first.setting-switch-what-bit-does-3}
The next result says that this is not one way of putting directions on a bit among many.

\begin{proposition}[Planes of directions on a two-state space]\label{c:synthesis.switch-planes-on-a-two-state-space}
\changedmark{BurntOrange}{changed since the deposit}
Let $F$ be a field in which $2 \ne 0$, $E = \mathrm{End}_{F}(F^{2})$, $E_{0} \subset E$ the maps of trace zero, and let a \emph{switch} be a pair $(a, b)$ in $E$ with $a^{2} = b^{2} = 0$, $ab + ba = 1$ (as in Theorem \ref{c:synthesis.set-and-clear-of-a-bit-generate-the-hyperbolic-plane}).
\begin{enumerate}
\item[(i)] If $V \subseteq E$ is a subspace of dimension at least $2$ in which every element squares to a scalar, then $V \subseteq E_{0}$. On $E_{0}$, $x^{2} = -\det x =: Q(x)$, and in the basis $h = wv - vw$, $v + w$, $v - w$ of a fixed switch $(v,w)$, $Q = \alpha^{2} + \sigma^{2} - \tau^{2}$.
\item[(ii)] Let $V \subseteq E_{0}$ be a plane. Its nonzero elements of square zero form two lines exactly when $V$ is spanned by a switch, one line exactly when $Q|_{V}$ is degenerate, and no line otherwise. No plane of $E_{0}$ is totally null.
\item[(iii)] If $F$ is ordered, an orthogonal basis of $V$ has squares of signs $(+,+)$, $(+,-)$ or $(+,0)$, and no other. If moreover every positive element of $F$ is a square, $V$ is spanned by a switch exactly in the case $(+,-)$.
\item[(iv)] $V$ generates $E$ as an algebra when $Q|_{V}$ is nondegenerate, and a subalgebra of dimension $3$ (conjugate to the upper triangular matrices) when it is degenerate.
\item[(v)] $V \mapsto \{\ker a, \ker b\}$, for any switch $(a,b)$ spanning $V$, is a well-defined bijection from the planes spanned by a switch to the unordered pairs of distinct lines of $F^{2}$, and any two such planes are conjugate under $\mathrm{GL}_{2}(F)$.
\end{enumerate}
\end{proposition}
\begin{proof}
Throughout, $(v, w)$ is a fixed switch in $E$, for instance set and clear in a basis of $F^{2}$ (Theorem \ref{c:synthesis.set-and-clear-of-a-bit-generate-the-hyperbolic-plane}(i)), $h = wv - vw$, and $hw = w$, $wh = -w$, $hv = -v$, $vh = v$, $h^{2} = 1$ by the relations of a switch.

\emph{(i)} Every $x \in E$ satisfies $x^{2} = (\operatorname{tr}x)\,x - \det x$, so $x^{2}$ is a scalar exactly when $x$ is a scalar or $\operatorname{tr}x = 0$. Suppose $x \in V$ has $\operatorname{tr}x \ne 0$. Then $x$ is a scalar $\lambda \ne 0$. Take $y \in V$ independent of $x$; $y$ is not a scalar, so $\operatorname{tr}y = 0$, and $x + y$ is not a scalar and has trace $2\lambda \ne 0$, so $(x + y)^{2}$ is not a scalar, a contradiction. On $E_{0}$, $x^{2} = -\det x$. By Theorem \ref{c:synthesis.set-and-clear-of-a-bit-generate-the-hyperbolic-plane}(ii) and (iv), $h$, $v + w$, $v - w$ is a basis of $E_{0}$ of pairwise anticommuting elements of squares $1$, $1$, $-1$, so $Q(\alpha h + \sigma(v + w) + \tau(v - w)) = \alpha^{2} + \sigma^{2} - \tau^{2}$.

\emph{(ii)} A nonzero $n \in E_{0}$ with $n^{2} = 0$ is a null vector of $Q$, and for $x, y \in E_{0}$, $xy + yx = (x + y)^{2} - x^{2} - y^{2}$ is the scalar $2B(x,y)$, $B$ the polar form of $Q$. Let $n_{1}, n_{2} \in E_{0}$ be independent with $n_{1}^{2} = n_{2}^{2} = 0$. A nonzero map of square zero on $F^{2}$ is conjugate to $w$, so choose the basis of $F^{2}$ with $n_{1} = w$. For $y = \alpha h + \beta w + \gamma v$, $wy + yw = \alpha(wh + hw) + \gamma(wv + vw) = \gamma$. If $n_{1}n_{2} + n_{2}n_{1} = 0$, then $n_{2} = \alpha h + \beta w$ and $n_{2}^{2} = \alpha^{2} + \alpha\beta(hw + wh) = \alpha^{2} = 0$, so $n_{2}$ is a multiple of $w$, which is excluded. Hence $n_{1}n_{2} + n_{2}n_{1} = \gamma \ne 0$. In particular $(n_{1} + n_{2})^{2} = \gamma \ne 0$: no plane is totally null. And $(n_{1}, \gamma^{-1}n_{2})$ is a switch spanning the plane. Conversely, on a plane spanned by a switch $(a, b)$, $Q(\alpha a + \beta b) = \alpha\beta$, with exactly the two null lines $Fa$ and $Fb$. If $Q|_{V}$ is degenerate, its radical is a line (not $V$, as $V$ is not totally null) of null vectors, and it is the only null line, since a second would be orthogonal to it and make $V$ totally null. If $Q|_{V}$ is nondegenerate and has a null vector $n$, choose $y \in V$ with $ny + yn \ne 0$; then $y - \frac{Q(y)}{ny + yn}\,n$ is null (expand the square) and independent of $n$. So a nondegenerate plane has no null line or two.

\emph{(iii)} Let $x, y$ be an orthogonal basis of $V$ (it exists since $2 \ne 0$). The plane $P = \operatorname{span}\{h, v + w\}$ has $Q > 0$ on its nonzero vectors, a sum of two squares in an ordered field. $V \cap P \ne 0$ by dimensions, so $Q$ takes a positive value on $V$, and not both of $Q(x), Q(y)$ are $\le 0$. If both were negative, $Q$ would be negative on the nonzero vectors of $V$, against $V \cap P \ne 0$. If one were negative and the other zero, or both zero, the same contradiction or (ii) applies ($Q \le 0$ on $V$, respectively $V$ totally null). This leaves $(+,+)$, $(+,-)$, $(+,0)$. A switch plane contains $a + b$ and $a - b$ of squares $1$ and $-1$, orthogonal: case $(+,-)$. Conversely if $Q(x) = \lambda > 0 > -\mu = Q(y)$ and $\mu/\lambda = \rho^{2}$, then $\rho x \pm y$ are independent null vectors, and (ii) applies.

\emph{(iv)} If $Q|_{V}$ is nondegenerate, take an orthogonal basis $x, y$: they anticommute and have nonzero scalar squares, so both are invertible. If $\alpha + \beta x + \gamma y + \delta xy = 0$, conjugating by $x$ changes the sign of the last two terms, so (as $2 \ne 0$) $\alpha + \beta x = 0$ and $(\gamma + \delta x)y = 0$, hence $\gamma + \delta x = 0$; since $x$ is not a scalar, all four coefficients vanish. So $1, x, y, xy$ are independent and span $E$. If $Q|_{V}$ is degenerate, take $n$ in the radical, a basis of $F^{2}$ with $n = w$, and $y \in V$ independent of $n$; by the computation in (ii), $y = \alpha h + \beta w$ with $\alpha \ne 0$, and $yw = \alpha w$, $wy = -\alpha w$, so the algebra generated is $\operatorname{span}\{1, w, h\}$, the upper triangular matrices in the basis $(s_{\mathrm{off}}, s_{\mathrm{on}})$.

\emph{(v)} By (ii) any switch $(a, b)$ spanning $V$ has $a$ and $b$ on the two null lines of $V$, so $\{\ker a, \ker b\}$ depends only on $V$. A nonzero map of square zero on $F^{2}$ has image equal to its kernel, a line. If $\ker a = \ker b = \ell$ then $\operatorname{im}a = \operatorname{im}b = \ell$ and $ab = ba = 0$, against $ab + ba = 1$: the two lines differ. A nonzero map of square zero is determined up to a scalar by its kernel, so the pair of lines determines the two null lines of $V$ and hence $V$: the map is injective. Given distinct lines $\ell_{0}, \ell_{1}$, choose $s_{\mathrm{off}} \in \ell_{0}$, $s_{\mathrm{on}} \in \ell_{1}$; set and clear in this basis are a switch with $\ker w = \ell_{0}$, $\ker v = \ell_{1}$: the map is onto. $\mathrm{GL}_{2}(F)$ acts transitively on pairs of distinct lines, and conjugation carries switch planes to switch planes compatibly with the map.

\end{proof}

\phantomsection\label{c:note.geometry-first.setting-switch-what-bit-does-4}
\noindent\emph{How the proof goes.} Every $2 \times 2$ matrix satisfies $x^{2} = (\operatorname{tr} x)\,x - \det x$, so a square that is a scalar forces trace zero; the rest is a computation on the maps of trace zero. The proof is in Appendix~\ref{prf:prop:planes}.

So on a system with two states, over the real numbers, a plane of directions that generates all the operations is either of signature $(2,0)$, which contains no operation of square zero, or it is a hyperbolic plane spanned by a switch, and the hyperbolic planes are exactly the choices of what the two states are. The pair of directions with squares $+1$ and $-1$ is the pair that goes with a bit that can be set and cleared.

\begin{theorem}[The integer Clifford algebras contain no idempotent and no switch]\label{c:synthesis.integer-clifford-algebras-contain-no-switch}
\changedmark{BurntOrange}{changed since the deposit}
For a commutative ring $R$ and $n = p + q \ge 1$ let $\Cl_{R}(p,q)$ be the free $R$-module on the blades $e_{S}$, $S \subseteq \{1, \dots, n\}$, with the product $e_{S}e_{T} = (-1)^{f(S,T)}e_{S \triangle T}$, $f(S,T) = \#\{(i,j) \in S \times T : i > j\} + |S \cap T \cap N|$, $N = \{p + 1, \dots, n\}$ (the algebra of Theorem \ref{c:synthesis.division-gives-the-rationals}). A \emph{switch} is a pair $(a, b)$ with $a^{2} = b^{2} = 0$, $ab + ba = 1$.
\begin{enumerate}
\item[(i)] $\Cl_{\Z}(p,q)$ has no idempotent other than $0$ and $1$.
\item[(ii)] If $ab + ba = 1$ for some $a, b \in \Cl_{R}(p,q)$, then $2$ is a unit of $R$. In particular $\Cl_{\Z}(p,q)$ contains no switch.
\item[(iii)] If $2$ is a unit of $R$ and $p, q \ge 1$, then $\tfrac12(e_{1} + e_{p+1})$, $\tfrac12(e_{1} - e_{p+1})$ is a switch.
\item[(iv)] Let $M$ be the free $\Z$-module on two states with set $v$ and clear $w$. The homomorphism $\Cl_{\Z}(1,1) \to \mathrm{End}_{\Z}(M)$ with $e_{1} \mapsto v + w$, $e_{2} \mapsto v - w$ is injective, and its image is the integer matrices whose diagonal entries have the same parity and whose off-diagonal entries have the same parity, a subring of index $4$.
\end{enumerate}
\end{theorem}
\begin{proof}
Write $\bar{A} = \Cl_{R}(p,q)/2\Cl_{R}(p,q)$. It is free over $R/2R$ on the $e_{S}$ with $e_{S}e_{T} = e_{S \triangle T}$, since $-1 = 1$ in $R/2R$: a commutative ring in which every $e_{S}$ squares to $1$.

\emph{(ii)} In the commutative ring $\bar{A}$ the image of $ab + ba$ is $2\bar{a}\bar{b} = 0$. So $1 \in 2\Cl_{R}(p,q)$, and the coefficient of $e_{\emptyset}$ in the free module gives $r \in R$ with $2r = 1$. For $R = \Z$ there is no such $r$.

\emph{(i)} Take $R = \Z$ and let $x = \sum_{S} x_{S}e_{S}$ be idempotent. $\bar{A}$ is commutative of characteristic $2$, so squaring is additive, and
\[
\bar{x}^{2} = \sum_{S} \bar{x}_{S}^{2}\,e_{S}^{2} = \Bigl(\sum_{S} \bar{x}_{S}\Bigr)\cdot 1 .
\]
Hence $\bar{x} = \bar{x}^{2}$ is $0$ or $1$. If $\bar{x} = 0$, then $x = 2y$ and $x = x^{k} = 2^{k}y^{k}$ for every $k \ge 1$, so every integer coefficient of $x$ is divisible by every power of $2$ and $x = 0$. If $\bar{x} = 1$, the idempotent $1 - x$ reduces to $0$, so $x = 1$.

\emph{(iii)} $e_{1}^{2} = 1$, $e_{p+1}^{2} = -1$ and $e_{1}e_{p+1} = -e_{p+1}e_{1}$, so $\bigl(\tfrac12(e_{1} \pm e_{p+1})\bigr)^{2} = \tfrac14(1 - 1) = 0$ and the anticommutator of the two is $\tfrac14\bigl(2e_{1}^{2} - 2e_{p+1}^{2}\bigr) = 1$.

\emph{(iv)} $e_{+} = v + w$ and $e_{-} = v - w$ anticommute with squares $1$ and $-1$ (Theorem \ref{c:synthesis.set-and-clear-of-a-bit-generate-the-hyperbolic-plane}(iv), recomputed here: $(v \pm w)^{2} = \pm(vw + wv) = \pm 1$, $(v + w)(v - w) = wv - vw = -(v - w)(v + w)$), so $e_{\emptyset} \mapsto 1$, $e_{1} \mapsto e_{+}$, $e_{2} \mapsto e_{-}$, $e_{12} \mapsto e_{+}e_{-}$ respects the blade product and is a ring homomorphism. In the basis $(s_{\mathrm{off}}, s_{\mathrm{on}})$ the image of $\alpha + \beta e_{1} + \gamma e_{2} + \delta e_{12}$ has diagonal entries $\alpha + \delta$, $\alpha - \delta$ and off-diagonal entries $\beta - \gamma$ (row off, column on) and $\beta + \gamma$. The map $(\alpha, \beta, \gamma, \delta) \mapsto (\alpha + \delta, \alpha - \delta, \beta - \gamma, \beta + \gamma)$ is injective on $\Z^{4}$ with determinant of absolute value $4$, and its image is exactly the quadruples whose first two entries have the same parity and whose last two have the same parity.

\end{proof}

\phantomsection\label{c:note.geometry-first.setting-switch-what-bit-does-5}
\noindent\emph{How the proof goes.} Modulo $2$ the algebra is commutative and every blade squares to $1$. So there the square of an element is a scalar, the sum of its coefficients, and an idempotent is $0$ or $1$ modulo $2$, which a $2$-adic argument lifts to (i); and $ab + ba = 2ab$ vanishes, so $ab + ba = 1$ is possible only if $2$ is invertible, which is (ii). Part (iii) is a two-line computation. Example~\ref{c:note.geometry-first.ladder-needs-halves-set-bit} shows the case $p = q = 1$ by hand. The proof is in Appendix~\ref{prf:thm:nobit}.

\begin{example}[The ladder needs halves to set a bit]\label{c:note.geometry-first.ladder-needs-halves-set-bit}
Send the two directions of $\Cl_{\Z}(1,1)$ to $e_{+}$ and $e_{-}$ of Example~\ref{c:note.geometry-first.switch-two-two-tables}. The element $\alpha + \beta e_{1} + \gamma e_{2} + \delta e_{1}e_{2}$, with integer coefficients, becomes
\[
\begin{pmatrix} \alpha + \delta & \beta - \gamma \\ \beta + \gamma & \alpha - \delta \end{pmatrix} .
\]
The exchange $e_{+}$ is there, with $\beta = 1$. Setting the bit needs $\beta + \gamma = 1$ and $\beta - \gamma = 0$, so $\beta = \gamma = \tfrac12$; the test for on needs $\alpha = \tfrac12$ and $\delta = -\tfrac12$. Over the integers the ladder can exchange the two states but can neither set the bit nor test it. With halves it can do both, and $\tfrac12(e_{1} + e_{2})$, $\tfrac12(e_{1} - e_{2})$ is the switch of Theorem~\ref{c:synthesis.integer-clifford-algebras-contain-no-switch}(iii).
\end{example}

\begin{remark}[Two orders of construction: the ladder and the switch]\label{c:note.synthesis.two-orders-of-construction-ladder-and-switch}
\changedmark{BurntOrange}{changed since the deposit}
The foundation's ladder takes the bit (F1), then one bit of non-commutativity (F2), then superposition over the integers (F3), and reaches $\Cl_{\Z}(p,q)$. Three results filed with this note separate that order of construction from a second one.

\emph{The switch.} Acting on a two-state module by set $v$ and clear $w$ gives $v^{2} = w^{2} = 0$, $vw + wv = 1$, the full endomorphism ring over every commutative ring, and the two directions $v \pm w$ of squares $\pm 1$, with no fraction (Theorem \ref{c:synthesis.set-and-clear-of-a-bit-generate-the-hyperbolic-plane}). On two states, the planes of directions containing two independent operations of square zero are exactly the planes of such pairs, one for each unordered pair of states (Theorem \ref{c:synthesis.switch-planes-on-a-two-state-space}).

\emph{The ladder contains no switch over the integers.} $\Cl_{\Z}(p,q)$ has no idempotent other than $0$, $1$, and $ab + ba = 1$ in $\Cl_{R}(p,q)$ forces $2$ to be a unit of $R$ (Theorem \ref{c:synthesis.integer-clifford-algebras-contain-no-switch}). So in the order F1, F2, F3 the direction is the reversible operation (undone, up to the sign, by repeating it), the algebra can exchange the two states, and it contains no test of which state holds and no set or clear until $2$ is inverted, which the step to $\Q$ does. In the other order the tests $vw$, $wv$ and the two irreversible operations are present at once, over the integers, and the reversible directions are their sum and difference. For one switch the first algebra is a subring of index $4$ in the second.

\emph{What is open.} The two orders agree over every ring containing $\tfrac12$. Which the foundation should take is not decided: F2 and F3 stand as filed, and the switch is recorded beside them. A foundation rooted on the switch would replace ``one central sign bit between distinct directions'' by ``each distinction can be set and cleared'', and would have to recover the sign and the anticommutation of distinct switches (the Jordan and Wigner sign), which a single switch does not supply.

\emph{Status.} That the on/off reading is physically significant (the exclusion principle as ``cannot be set twice'') is a reading, argued and not proved. The three results are algebra; none uses a rule of transport and none is evidence for the transport postulate.
\end{remark}

\section{What the tests leave of a switch}

\phantomsection\label{c:note.geometry-first.setting-what-tests-leave-switch}
Remark \ref{c:note.synthesis.two-orders-of-construction-ladder-and-switch} set the ladder and the switch side by side. One map joins them. Keep the tests of $n$ switches and discard everything that changes a state. What is left is the algebra of functions on the configurations of premise (F1), and when $2 \neq 0$ it carries their comparison.

\phantomsection\label{c:note.geometry-first.setting-what-tests-leave-switch-2}
For $n = 1$ this is the bit of Definition \ref{def:switch}: $P_{1} = vw$ is the test for on, and $\omega_{1} = wv - vw$ is the element $e_{+}e_{-}$ of Theorem \ref{c:synthesis.set-and-clear-of-a-bit-generate-the-hyperbolic-plane}(iv).

\phantomsection\label{c:note.geometry-first.setting-what-tests-leave-switch-3}
\noindent\emph{How the proof goes.} The tests are diagonal in the basis of configurations and the joint tests project onto its lines, so the pinching keeps the diagonal part; the rest is bookkeeping with signs. Example~\ref{ex:pinch} is one switch. The proof is in Appendix~\ref{prf:prop:pinching}.

So what the pinching leaves of $n$ switches is the commutative algebra $D$ of $R$-valued functions on the configurations of (F1). The joint tests, one for each configuration, are a basis of it. When $2 \neq 0$ in $R$ the parities carry the comparison of configurations faithfully, and $\omega_{T}^{2} = 1$ is the rule that a distinction made twice returns to nothing; when $2 = 0$ every parity is $1$ and $D$ does not see the comparison.

\phantomsection\label{c:note.geometry-first.setting-what-tests-leave-switch-4}
$D$ is not (F1) itself. It holds more: the scalars of $R$, and with them a sign, where the bit of (F1) has neither. And the comparison is read from the parities, which are built from the tests $P_{i}$; it is not determined by $D$ as an algebra alone. What $D$ does not hold is any operation on a state. It is commutative, so the order of two distinctions does not matter, which is what (F2) adds; and by (iii) each of its elements maps every $R\,s_{S}$ into itself, so nothing in it sets, clears or exchanges a state. It can be read and cannot be written. To call $D$ the bit of (F1), seen with scalars, is a reading and is used below as one. In the Clifford algebra, by (v), $D$ is spanned by the products of the bivectors $e_{i}e_{n+i}$ of the hyperbolic planes.

\begin{remark}[The two orders, joined]\label{c:note.geometry-first.two-orders-joined}
The two orders of Remark \ref{c:note.synthesis.two-orders-of-construction-ladder-and-switch} meet in the operations on a register, and $E$ tells them apart. The switch order starts from the $v_{i}$ and $w_{i}$. The ladder starts at (F1), which on the reading above is the image of $E$ without its scalars, and its premises (F2) and (F3) add, for each distinction, one combination of set and clear: the direction $L_{i} = v_{i} + \varepsilon_{i}w_{i}$ of part (vi), with the signature as the sign $\varepsilon_{i}$. The other combination $L_{i}'$, of the opposite square, is not a left multiplication, unless $2 = 0$ in $R$, when it coincides with $L_{i}$: if it were left multiplication by $y$ then $y = L_{i}'e_{\emptyset} = e_{i}$, and $L_{i} - L_{i}' = 2\varepsilon_{i}w_{i}$.

Nor does the ladder's algebra contain any test of its own distinctions, over any ring: by (vi) the pinching sends every left multiplication to its scalar part, so a left multiplication that commutes with every test is a scalar. (The tests that Theorem \ref{c:synthesis.integer-clifford-algebras-contain-no-switch}(iii) supplies when $2$ is invertible belong to a different bit; they are not diagonal on the configurations.) Which end the foundation should take is still a choice, and (F1) to (F4) stand as stated.

The map is already in the series. For one switch,
\[
E(x) \;=\; PxP + (1-P)\,x\,(1-P)
\]
is the descent of T5 \cite{ButlerT5} (Proposition 5.5 there), and for $n$ switches it is T5's multi-mode descent (Definition 9.1 there), whose range T5 calls the classical algebra of occupation records. Read with Proposition \ref{prop:pinching}: the descent of a register returns the functions on the configurations of (F1), the bare bit on the reading above, and the ascent, which adjoins the pairs $(v_{i}, w_{i})$, supplies its operations. This is a comparison of definitions. T5's use of the descent, on the algebras and states of a black hole, is not treated here, and Proposition \ref{prop:pinching} does not depend on it.

\emph{Status.} Proposition \ref{prop:pinching} is proved, and every step of it is elementary. That it matters for physics is a reading. Nothing here uses a rule of transport, and none of it is evidence for the transport principle.
\end{remark}

\section{Position in the series}

\phantomsection\label{c:note.geometry-first.setting-position-series}
This paper is T0a of the series, a companion note to the primer T0 \cite{ButlerT0}. It was written as the primer's closing section and is kept apart so that the primer stays what it was meant to be, an expository start with nothing new in it.

\emph{This note relies on:} T0 \cite{ButlerT0} for its language and notation, cited where used. No result of the series enters a proof here.

\emph{Papers that rely on this one:} none, in their proofs. T1's Example 3.4 \cite{ButlerT1} points here for the reading of the hyperbolic pair as the two operations of a bit, and Remark \ref{c:note.geometry-first.two-orders-joined} points to T5 \cite{ButlerT5}, whose descent is the pinching of Proposition \ref{prop:pinching}; the proposition does not depend on T5.

\phantomsection\label{c:note.geometry-first.setting-position-series-2}
\textbf{Where the series uses this.} The switch of Theorem \ref{c:synthesis.set-and-clear-of-a-bit-generate-the-hyperbolic-plane} is the hyperbolic pair that T1 \cite{ButlerT1} adjoins to repair a failed division. In T1's Example 3.4 the product $wv$, here the test for off, is the idempotent $e$ with $ev = 0$ and $v(1 - e) = 0$; that example and this note point to each other. The failure of division at $v$ and $w$, and its repair, remain T1's subject and are not treated here. T3 \cite{ButlerT3} iterates the adjunction, one switch per step. The pinching of Proposition \ref{prop:pinching} is the descent of T5 \cite{ButlerT5}. Theorem \ref{c:synthesis.integer-clifford-algebras-contain-no-switch} is arithmetic of the integer algebra, and nothing in this note uses a rule of transport: none of it is evidence for the transport principle.

\section{Machine verification}

\phantomsection\label{c:note.geometry-first.setting-machine-verification}
The computations of this note are checked by the script \texttt{verify\_T00a.py}, deposited alongside it, which imports the series' shared exact-arithmetic library \texttt{GA.py} and works in exact rational arithmetic, with no floating point anywhere. A model of premise (F2) is built for every signature with $n \leq 4$, which shows the premise is consistent, and the formulas of Theorem \ref{thm:bitsign} for squares and commutators are compared with it; the blade product and the one rule of Theorem \ref{c:synthesis.signed-directions-and-the-integer-clifford-algebra} are compared with \texttt{GA.py}. For the completion step the script checks the inequalities of the square-root argument and the independence of $(p,q,r)$ from the order of rejection (Theorem \ref{c:synthesis.completion-recovers-the-real-clifford-algebras}), the reduction of Lemma \ref{lem:bitpoly}(b) on known examples, the square-root formula and the real coefficients of $g_{t}$ (Theorem \ref{thm:bitroots}), and the relations used in Theorem \ref{c:synthesis.frobenius-and-universality-in-clifford-algebras}. The identities of Theorem \ref{c:synthesis.set-and-clear-of-a-bit-generate-the-hyperbolic-plane}, the case analysis of Proposition \ref{c:synthesis.switch-planes-on-a-two-state-space} over the rationals, Theorem \ref{c:synthesis.integer-clifford-algebras-contain-no-switch} and the index of Remark \ref{c:note.synthesis.two-orders-of-construction-ladder-and-switch} are checked on finite ranges of integers, the mod-$2$ identity exhaustively for $n \leq 3$. Proposition \ref{prop:pinching} is checked on the plain and the signed register for $n \leq 3$ over the integers, including the index of part (iv), with the statements that depend on the ring (the injectivity in (iv) and the two claims of Remark \ref{c:note.geometry-first.two-orders-joined}) checked over $\Z/2$, $\Z/3$, $\Z/4$ and $\Z/6$ for $n \leq 2$; its part (v) in $\Cl(n,n)$ for $n \leq 3$ with \texttt{GA.py}, and its part (vi) against \texttt{GA.py} for every choice of signs with $n \leq 3$. These searches support the proofs and do not replace them; the limiting arguments that use completeness are proved, not computed. The claims about the switch and their conditions of failure were registered before any of them was run (\texttt{preregistration\_bit\_switch.md}, deposited with this note); Proposition \ref{prop:pinching} was not registered in advance. Not checked by machine, their proofs being short and structural: the universal properties of Theorems~\ref{c:synthesis.signed-directions-and-the-integer-clifford-algebra}(iii) and~\ref{c:synthesis.completion-recovers-the-real-clifford-algebras}(iv), Lemma~\ref{lem:bitpoly}(a), Theorem~\ref{thm:bitroots}(iv), and the trace argument for bosons. Fifty-three checks for the results. The worked examples are checked by the same script, each as printed: the table of Example~\emph{Two bits} [T0a, Ex A] against Theorem~\ref{thm:bitsign}(iii), the squares and inverse of Examples~\ref{c:note.geometry-first.one-rule-hand} and~\ref{c:note.geometry-first.vector-needs-half}, the roots and Hamilton's relations of Example~\ref{c:note.geometry-first.roots-cl-0-1-hamiltons-relations}, the tables of Examples~\ref{c:note.geometry-first.switch-two-two-tables} and~\ref{c:note.geometry-first.ladder-needs-halves-set-bit} (with the uniqueness of the halves), and the pinching of Example~\ref{ex:pinch}; eight checks. Sixty-one checks in all, all passing.

\section{Proofs}

\phantomsection\label{c:note.geometry-first.setting-proofs}
This appendix collects the proofs of the results stated in the section ``Order, and the arrival of the sign'' to~the section ``What the tests leave of a switch'', in the order of the text.

\section{Results used from elsewhere}
The proofs above use the following results of the series, stated briefly; each is proved in the paper named.

\begin{fact}[Minimal nondegenerate completion]\label{c:conformal.minimal-nondegenerate-completion}
A real quadratic space with radical of dimension $r$ and nondegenerate part of signature $(p,q)$ has minimal nondegenerate completions of exactly $r$ extra dimensions, all of signature $(p+r,q+r)$ and unique up to an isometry fixing the space, so $\Cl(p,q,r)$ embeds in $\Cl(p+r,q+r)$, minimally among completion-induced embeddings.
\emph{Source:} T2, Thm 2.3 (\href{https://doi.org/10.5281/zenodo.22714886}{doi:10.5281/zenodo.22714886}).
\end{fact}

\begin{fact}[The blade basis, degenerate forms included]\label{c:std-algebra.blade-basis}
For any basis of $V$ the $2^{n}$ increasing products $e_I$ form a basis of $\Cl(V, Q)$, degenerate forms included; so $V \hookrightarrow \Cl(V, Q)$.
\emph{Source:} the algebra core of the Physics Graph, proved there with an exact check; standard, see C. Chevalley, The Algebraic Theory of Spinors, Columbia University Press (1954), Ch. II (the Clifford algebra of a quadratic space).
\end{fact}

\begin{fact}[The complex numbers inside Clifford algebras]\label{c:std-algebra.complex-numbers}
$\Cl(0,1) \cong \C$ and $\Cl^{+}(2,0) = \operatorname{span}\{1, e_1 e_2\} \cong \C$, the unit bivector playing the imaginary unit.
\emph{Source:} the algebra core of the Physics Graph, proved there with an exact check; standard, see C. Doran and A. Lasenby, Geometric Algebra for Physicists, Cambridge University Press (2003), Ch. 2 (geometric algebra in two and three dimensions).
\end{fact}

\begin{fact}[The Clifford algebra exists and is unique]\label{c:std-algebra.existence-uniqueness}
$\Cl(V, Q) = T(V)/\langle v \otimes v - Q(v) \rangle$ exists and is unique up to unique isomorphism; the wedge-plus-contraction operators on $\Lambda V$ give a concrete representation.
\emph{Source:} the algebra core of the Physics Graph, proved there with an exact check; standard, see C. Chevalley, The Algebraic Theory of Spinors, Columbia University Press (1954), Ch. II (the Clifford algebra of a quadratic space).
\end{fact}

\begin{fact}[The Fock representation, and the Fock space as a minimal left ideal]\label{c:std-algebra.fock-representation}
$\Cl(\R^{k,k}) \cong \End(\Lambda W)$ by wedge and contraction, and the Fock space is the minimal left ideal $\Cl\, f$ of the vacuum idempotent $f = \prod_i v_i w_i$.
\emph{Source:} the algebra core of the Physics Graph, proved there with an exact check; standard, see C. Chevalley, The Algebraic Theory of Spinors, Columbia University Press (1954), Ch. II and III (spinors, pure spinors).
\end{fact}

\begin{fact}[Isometries extend to algebra homomorphisms]\label{c:std-algebra.functoriality}
An isometry $f$ extends uniquely to an algebra homomorphism $\Cl(f)$, injective when $f$ is; $\mathrm{O}(V, Q)$ acts by automorphisms.
\emph{Source:} the algebra core of the Physics Graph, proved there with an exact check; standard, see P. Lounesto, Clifford Algebras and Spinors, 2nd ed., London Mathematical Society Lecture Note Series 286, Cambridge University Press (2001), Ch. 14 (definitions of the Clifford algebra).
\end{fact}

\begin{fact}[Matrix units give a matrix algebra over the corner]\label{c:std-algebra.matrix-units}
A full $m \times m$ system of matrix units $E_{ij}$ in $A$ gives $A \cong \Mat_m(E_{11} A E_{11})$ via $b \mapsto (E_{1i} b E_{j1})$; e.g. $\Cl(1,1) \cong \Mat_2(\R)$.
\emph{Source:} the algebra core of the Physics Graph, proved there with an exact check; standard, see T. Y. Lam, Lectures on Modules and Rings, Graduate Texts in Mathematics 189, Springer, New York (1999), matrix rings, idempotents and Morita theory.
\end{fact}

\begin{fact}[Orthogonal vectors anticommute]\label{c:std-algebra.orthogonal-anticommute}
$uv = -vu$ iff $B(u,v) = 0$; in an orthogonal basis $e_I e_J = \pm \prod_{I \cap J} Q(e_i)\; e_{I \triangle J}$ with the sorting sign.
\emph{Source:} the algebra core of the Physics Graph, proved there with an exact check; standard, see C. Doran and A. Lasenby, Geometric Algebra for Physicists, Cambridge University Press (2003), Ch. 4 (foundations of geometric algebra).
\end{fact}

\begin{fact}[Polarisation]\label{c:std-algebra.polarisation}
The polar form is symmetric with $B(v,v) = Q(v)$, every symmetric bilinear form arises this way, and in a basis $Q(x) = x^{\mathsf T} G x$.
\emph{Source:} the algebra core of the Physics Graph, proved there with an exact check; standard, see T. Y. Lam, Introduction to Quadratic Forms over Fields, Graduate Studies in Mathematics 67, American Mathematical Society (2005), Ch. I and II (quadratic forms; the real case).
\end{fact}

\begin{fact}[The quaternions inside Clifford algebras]\label{c:std-algebra.quaternions}
$\Cl^{+}(3,0) = \operatorname{span}\{1, e_3e_2, e_1e_3, e_2e_1\} \cong \HH$ and $\Cl(0,2) \cong \HH$, by Hamilton's relations.
\emph{Source:} the algebra core of the Physics Graph, proved there with an exact check; standard, see C. Doran and A. Lasenby, Geometric Algebra for Physicists, Cambridge University Press (2003), Ch. 2 (geometric algebra in two and three dimensions).
\end{fact}

\begin{fact}[Null pairs satisfy the canonical anticommutation relations]\label{c:std-algebra.witt-basis-car}
For a Witt basis ($B(v_i,v_j) = B(w_i,w_j) = 0$, $B(v_i,w_j) = \tfrac12\delta_{ij}$), $\{v_i,v_j\} = \{w_i,w_j\} = 0$ and $\{v_i,w_j\} = \delta_{ij}$ in $\Cl(V,Q)$.
\emph{Source:} the algebra core of the Physics Graph, proved there with an exact check; standard, see C. Doran and A. Lasenby, Geometric Algebra for Physicists, Cambridge University Press (2003), Ch. 9 (multiparticle states).
\end{fact}

\begin{fact}[The exceptional division algebras]\label{c:transport.exceptional-division-algebras}
The only real Clifford algebras without zero divisors are $\R$, $\C$ and $\HH$, and the only ones without nilpotents are these together with $\R \oplus \R$ and $\HH \oplus \HH$.
\emph{Source:} T1, Prop 3.2 (\href{https://doi.org/10.5281/zenodo.21743863}{doi:10.5281/zenodo.21743863}).
\end{fact}

\section*{Changes since the deposited version}
\addcontentsline{toc}{section}{Changes since the deposited version}
Since the deposited version 1.0: 13 results changed, 12 cited results added and 1 deposited result withdrawn. The comparison is with the source of version 1.0 on Zenodo, \href{https://doi.org/10.5281/zenodo.24402873}{doi:10.5281/zenodo.24402873}. A result counts as revised when the words or mathematics of its statement or proof differ from that source; a citation added in a parenthesis does not count. A result the paper cites from elsewhere is compared with the deposit of the paper that proves it. Each revision is classed as corrected, changed or reworded; the classes were assigned by a review of each difference and the store's record of why it was made, and a revision not yet reviewed is counted as changed.
\subsection*{Changed}
The content changed by a later decision, not because the deposited version was wrong.
\begin{itemize}\setlength{\itemsep}{2pt}
\item Postulate~\ref{c:synthesis.continuum-completion-premise} (\emph{Premise (F4): the coefficients are completed to a complete ordered field}), Premise (F4) in the deposit: statement.
\item Postulate~\ref{c:synthesis.premise-minimal-non-commutativity} (\emph{Premise (F2): minimal non-commutativity}), Premise (F2) in the deposit: statement.
\item Postulate~\ref{c:synthesis.premise-superposition-over-the-integers} (\emph{Premise (F3): superposition over the integers}), Premise (F3) in the deposit: statement.
\item Postulate~\ref{c:synthesis.premise-the-bit} (\emph{Premise (F1): the bit}), Premise (F1) in the deposit: statement.
\item Theorem~\ref{c:synthesis.signed-directions-and-the-integer-clifford-algebra} (\emph{The group of signed directions and its integer Clifford algebra}), Thm 4.1 in the deposit: statement and proof.
\item Proposition~\ref{c:synthesis.division-gives-the-rationals} (\emph{Division in the integer Clifford algebra gives the rationals}), Prop 5.1 in the deposit: statement and proof.
\item Theorem~\ref{c:synthesis.completion-recovers-the-real-clifford-algebras} (\emph{Completion recovers the real Clifford algebras; Sylvester's law of inertia in GA}), Thm 6.1 in the deposit: statement and proof.
\item Theorem~\ref{c:synthesis.frobenius-and-universality-in-clifford-algebras} (\emph{Frobenius in GA; every finite-dimensional associative real algebra in a corner of Cl(k,k)}), Thm 6.4 in the deposit: statement and proof.
\item Example~\ref{c:note.geometry-first.roots-cl-0-1-hamiltons-relations} (\emph{Roots in $\Cl(0,1)$, and Hamilton's relations}), Ex D in the deposit: statement.
\item Theorem~\ref{c:synthesis.set-and-clear-of-a-bit-generate-the-hyperbolic-plane} (\emph{Set and clear of a bit generate the hyperbolic plane}), Thm 7.2 in the deposit: statement and proof.
\item Proposition~\ref{c:synthesis.switch-planes-on-a-two-state-space} (\emph{Planes of directions on a two-state space}), Prop 7.3 in the deposit: statement and proof.
\item Theorem~\ref{c:synthesis.integer-clifford-algebras-contain-no-switch} (\emph{The integer Clifford algebras contain no idempotent and no switch}), Thm 7.4 in the deposit: statement and proof.
\item Remark~\ref{c:note.synthesis.two-orders-of-construction-ladder-and-switch} (\emph{Two orders of construction: the ladder and the switch}), Rem 7.5 in the deposit: statement.
\end{itemize}
\noindent In addition, 12 results from elsewhere in the series or the algebra core are stated here that the deposited paper cited or used without stating. These are not marked.
\subsection*{Withdrawn}
\begin{itemize}
\item Example G (\emph{note.geometry-first.one-switch-pinched}): withdrawn from the store.
\end{itemize}
\begin{thebibliography}{99}
\bibitem{Galileo1623} \emph{Reference Galileo1623 is not in the store's registry.}
\bibitem{Wheeler1990} \emph{Reference Wheeler1990 is not in the store's registry.}
\bibitem{Hestenes2003Oersted} \emph{Reference Hestenes2003Oersted is not in the store's registry.}
\bibitem{ButlerT0} L.~Butler,
\emph{An Introduction to Geometric Algebra: A Primer for the Transport Papers},
Transport Papers T0 (2026).
\bibitem{SpencerBrown1969} \emph{Reference SpencerBrown1969 is not in the store's registry.}
\bibitem{Kauffman2021} \emph{Reference Kauffman2021 is not in the store's registry.}
\bibitem{Weizsacker1985} \emph{Reference Weizsacker1985 is not in the store's registry.}
\bibitem{AlbuquerqueMajid2002} \emph{Reference AlbuquerqueMajid2002 is not in the store's registry.}
\bibitem{Lam2005} T.~Y.~Lam,
\emph{Introduction to Quadratic Forms over Fields},
Graduate Studies in Mathematics 67, American Mathematical Society (2005).
\bibitem{Ebbinghaus1991} \emph{Reference Ebbinghaus1991 is not in the store's registry.}
\bibitem{Frobenius1878} \emph{Reference Frobenius1878 is not in the store's registry.}
\bibitem{Palais1968} \emph{Reference Palais1968 is not in the store's registry.}
\bibitem{ButlerT3} L.~Butler,
\emph{The Tower and the CAR Limit: Iterated Transport, the Frozen Clock, and Second Quantisation},
Transport Papers T3 (2026), doi:\href{https://doi.org/10.5281/zenodo.22721556}{10.5281/zenodo.22721556}.
\bibitem{JordanWigner1928} \emph{Reference JordanWigner1928 is not in the store's registry.}
\bibitem{Bass1974} \emph{Reference Bass1974 is not in the store's registry.}
\bibitem{Knus1991} \emph{Reference Knus1991 is not in the store's registry.}
\bibitem{ButlerT5} L.~Butler,
\emph{Division by Zero at the Gravitational Singularity: The Transport Principle Applied to the Black Hole, and the Information Paradox as a Category Error},
Transport Papers T5 (2026), doi:\href{https://doi.org/10.5281/zenodo.22721790}{10.5281/zenodo.22721790}.
\bibitem{ButlerT1} L.~Butler,
\emph{The Transport Principle: Total Division in Clifford Algebras via Minimal Completion},
Transport Papers T1 (2026).
\end{thebibliography}
\end{document}